\documentclass[11pt]{article}
\usepackage[T1]{fontenc}
\usepackage{lmodern}
\usepackage[margin=1in]{geometry}
\usepackage{amsmath,amssymb,amsthm}
\usepackage{graphicx,xcolor,microtype,xurl,pict2e}
\usepackage[hidelinks]{hyperref}
\hypersetup{pdftitle={Universal-cover splitting for compact Kahler manifolds},
  pdfauthor={OpenAI},pdfsubject={Compatible products, integrable tangent splittings, and canonical contractions}}
\pdftrailerid{}
\pdfinfoomitdate=1
\pdfsuppressptexinfo=15
\newtheorem{theorem}{Theorem}[section]
\newtheorem{proposition}[theorem]{Proposition}
\newtheorem{lemma}[theorem]{Lemma}
\newtheorem{corollary}[theorem]{Corollary}
\theoremstyle{definition}
\newtheorem{definition}[theorem]{Definition}
\theoremstyle{remark}
\newtheorem{remark}[theorem]{Remark}
\newcommand{\C}{\mathbb C}
\newcommand{\R}{\mathbb R}
\newcommand{\D}{\mathbb D}
\newcommand{\Pp}{\mathbb P}
\newcommand{\Z}{\mathbb Z}
\newcommand{\dd}{\mathrm d}
\newcommand{\id}{\operatorname{id}}
\newcommand{\reg}{\mathrm{reg}}
\DeclareMathOperator{\Ric}{Ric}
\DeclareMathOperator{\tr}{tr}
\DeclareMathOperator{\rank}{rank}
\DeclareMathOperator{\diam}{diam}
\DeclareMathOperator{\supp}{supp}
\DeclareMathOperator{\pr}{pr}
\newcommand{\eps}{\varepsilon}
\newcommand{\OO}{\mathcal O}
\newcommand{\norm}[1]{\left\lVert#1\right\rVert}
\numberwithin{equation}{section}
\allowdisplaybreaks[2]
\setlength{\emergencystretch}{2em}
\setcounter{tocdepth}{1}
\title{Universal-cover splitting for compact K\"ahler manifolds}
\author{OpenAI}
\date{September 23, 2026}
\begin{document}
\maketitle
\begin{abstract}
We prove the two-summand form of Beauville's compatible splitting conjecture.
If the tangent bundle of a compact connected K\"ahler manifold decomposes into
two integrable holomorphic subbundles of positive rank, then its ordinary
universal cover admits a product decomposition whose factor tangent bundles
are the lifted specified summands.
\end{abstract}
\tableofcontents
\clearpage
\section{Introduction}\label{sec:intro}

If a complex manifold is a product, its holomorphic tangent bundle is the
direct sum of the tangent bundles of the factors. The converse starts
with a specified holomorphic decomposition
\[
T_X=E_1\oplus E_2
\]
and asks whether the infinitesimal factors come from a product on the
universal cover. Compatibility with the given subbundles is part of the
question: an abstract product decomposition need not realize the
specified embeddings of $E_1$ and $E_2$ into $T_X$.
A subbundle is \emph{integrable} if its local holomorphic sections are
closed under the Lie bracket. When both summands are integrable,
holomorphic Frobenius gives simultaneous local product coordinates.
The problem is to continue these local products globally.

This question is closely related to the distinction between a local
product structure and a complete product geometry. For a complete
Riemannian metric, parallel complementary orthogonal distributions lead
to the de Rham decomposition \cite{deRham52}. A holomorphic splitting supplies no
complete metric making the distributions parallel. Compactness makes
every smooth metric complete, but it does not turn the given foliations
into parallel distributions. The missing step is therefore a
continuation theorem for the specified local products.

\subsection*{The compatible splitting problem}

Beauville formulated the compatible
universal-cover conjecture for compact K\"ahler manifolds with
integrable partial sums of the tangent summands
\cite[Section~2.3]{Beauville00}, and proved major positive cases, including
K\"ahler--Einstein manifolds and compact K\"ahler surfaces
\cite[Theorems~A and~C]{Beauville00}. In the K\"ahler--Einstein case,
Beauville refines an earlier splitting result of Yau \cite{Yau93}, as
he explains in his introduction. For two summands, the partial-sum
condition means precisely that both are integrable. For this formulation,
see also
\cite[Conjecture~1.4]{Horing26}.

The distinction between a product and a compatible product is already
visible in Beauville's example on $A\times\Pp^1$, where $A$ is an
abelian surface \cite[Section~2.2]{Beauville00}. Take independent
translation-invariant vector fields $U,V$ on $A$ and vector fields
$S,T$ on $\Pp^1$ with $[S,T]\ne0$, viewed as fields on the product.
The rank-two subbundle spanned by $U+S,V+T$ complements $T_{\Pp^1}$,
but is not integrable: its two generators have the nonzero vertical
bracket $[S,T]$. Thus even an underlying product need not realize a
specified tangent splitting.

Several later results establish global continuation under additional
geometric hypotheses. Druel treats projective manifolds whose tangent
bundle is a sum of line bundles, with the relevant partial sums
integrable, and obtains automatic integrability in the minimal
projective case \cite[Theorems~1 and~2]{Druel00}. Brunella, Pereira, and Touzet treat the case of a
line summand with integrable complement on a compact K\"ahler manifold
\cite[Theorem~1.1(i)]{BPT06}. Their work also gives uniformization
results when the tangent bundle is a sum of line bundles. Their proof
uses leafwise Poincar\'e metrics and transverse invariant metrics to
control continuation. Theorem~\ref{thm:main} below permits both summands
to have arbitrary positive rank.

In higher rank, Pereira and Touzet obtain a compatible Euclidean factor
when one involutive direct summand is Hermitian flat
\cite[Theorem~B]{PT13}. Druel, Pereira, Pym, and Touzet show that a
regular foliation with numerically flat tangent bundle---meaning that
the bundle and its dual are nef---is Hermitian flat and admits an
integrable complement, again producing a compatible Euclidean factor
\cite[Theorem~1.1]{DPPT24}.
A different continuation mechanism applies when a foliation has a compact
leaf. Druel, Pereira, Pym, and Touzet prove compatible
universal-cover splitting when one foliation has a compact leaf with
finite holonomy \cite[Theorem~1.4]{DPPT22}. Their argument uses global
Reeb stability and a flat Ehresmann connection. Our proof replaces
these additional flatness or leaf hypotheses by extension estimates
for maps adapted to the two distributions.

In the projective setting, integrability and product decomposition are
distinct issues. H\"oring proves automatic integrability for split
tangent bundles on non-uniruled projective manifolds
\cite[Theorem~1.3]{Horing07}. His proof combines the positivity theorem
for cotangent quotients of Campana and Peternell
\cite[Theorems~0.1 and~1.8]{CP11} with Demailly's integrability theorem
for holomorphic forms valued in the inverse of a pseudoeffective line
bundle \cite[Main Theorem]{Demailly02}. His recent work proves algebraic
integrability for tangent summands on normal $\mathbf Q$-factorial
projective klt varieties of Fano type; its product corollary, which does
not require $\mathbf Q$-factoriality, gives compatible products after
finite quasi-\'etale covers
\cite[Theorem~1.1 and Corollary~1.2]{Horing26}.
The companion \cite[Theorem~1.1]{OpenAIRCSplitting2026} proves that,
on a smooth connected rationally connected projective complex manifold,
both summands of every specified holomorphic splitting
$T_X=E_1\oplus E_2$ into positive-rank subbundles are integrable.
These results address
the geometry that forces a splitting to define foliations. In our main
theorem integrability is an explicit hypothesis, while the ambient
manifold is only assumed to be compact K\"ahler.

\begin{theorem}\label{thm:main}
Let $X$ be a compact connected K\"ahler manifold of complex dimension
$n\geq2$, and let
\[
T_X=E_1\oplus E_2
\]
be a specified decomposition into integrable holomorphic subbundles of
positive ranks $r_1,r_2$. If $\pi:\widetilde X\to X$ is the ordinary
universal covering map, there are connected simply connected complex
manifolds $Y_1,Y_2$ and a biholomorphism
$\Phi:\widetilde X\to Y_1\times Y_2$ such that
\[
\dd\Phi(\pi^*E_i)=\pr_i^*T_{Y_i}\qquad(i=1,2),
\]
where $\dd\pi$ identifies $T_{\widetilde X}$ with $\pi^*T_X$.
In particular, $\dim_\C Y_i=r_i$.
\end{theorem}

Theorem~\ref{thm:main} proves the two-summand form of Beauville's
compatible splitting conjecture. Its factors are the universal covers
of the intrinsic leaves through a common point. The leaves and factors
can be noncompact. The conclusion is about the ordinary universal cover;
it does not assert that a finite cover of $X$ is a product. If a summand
has rank zero, the corresponding factor is a point, so the positive-rank
formulation removes only the trivial case.

\begin{corollary}[Non-uniruled projective manifolds]\label{cor:non-uniruled}
Let $X$ be a smooth connected non-uniruled projective complex manifold,
and let $T_X=E_1\oplus E_2$ be a specified holomorphic splitting into
subbundles of positive ranks. Then both summands are integrable and the
ordinary universal cover has a compatible product decomposition as in
Theorem~\ref{thm:main}.
\end{corollary}

\begin{proof}
H\"oring's theorem \cite[Theorem~1.3]{Horing07} makes both summands
integrable. Since a smooth projective complex manifold is compact
K\"ahler, Theorem~\ref{thm:main} applies.
\end{proof}

\subsection*{The nef-and-big application}

For a smooth projective manifold, a line bundle is \emph{nef} if it has
nonnegative degree on every curve, and it is \emph{big} if its sections
have maximal asymptotic growth. When $K_X$ is nef and big, its
pluricanonical sections define a birational contraction $f:X\to Z$ to
the canonical model. We call $X^0=f^{-1}(Z_{\reg})$, the inverse image
of the smooth locus of $Z$, the \emph{regular canonical locus}.

Corollary~\ref{cor:non-uniruled} already gives the product conclusion
in this case. Our metric argument proves an additional statement:
the canonical K\"ahler--Einstein connection preserves both supplied
summands on $X^0$. This parallelism is the input for the alternative
continuation method in the appendices. We first record the product
conclusion.

\begin{corollary}\label{thm:nef-big}
Let $X$ be a smooth connected projective complex manifold of dimension
$n\geq2$ with nef and big canonical bundle $K_X$. Every specified
holomorphic splitting $T_X=E_1\oplus E_2$ into positive-rank subbundles
is integrable and induces a compatible product decomposition of the
ordinary universal cover as in Theorem~\ref{thm:main}.
\end{corollary}

When $K_X$ is ample, the negative K\"ahler--Einstein metric and
Beauville's theorem already give this conclusion. For nef and big $K_X$,
we use the canonical K\"ahler--Einstein approximation on a smooth
crepant resolution, as developed by Song \cite{Song14}.
The integrated squared second fundamental forms of the supplied summands
are $O(\varepsilon)$ along the approximation. Their limit vanishes on
the regular canonical locus, where the limiting metric makes both
summands parallel. The holomorphic Frobenius obstructions then vanish
on all of $X$. This supplies the precise hypothesis needed to apply
Theorem~\ref{thm:main}.
The approximate K\"ahler--Einstein method belongs to the stability
theory developed by Enoki \cite{Enoki88}; the passage from vanishing
second fundamental forms to a parallel splitting is used explicitly
by Guenancia \cite[Section~3.4, Step~4]{Guenancia16}.

\subsection*{The analytic continuation argument}

Local product coordinates alone do not show that a path in one foliation
can be transported all the way along a path in the other. We first solve
that extension problem for holomorphic maps, and then pass to continuous
paths. A holomorphic \emph{box} in two scalar parameters is a map $f$
whose $s$-derivative lies in $E_1$ and whose $t$-derivative lies in $E_2$.
In a simultaneous product chart it has the separated form
$(f_1(s),f_2(t))$. We do not impose an immersion or full-rank condition.

The first analytic step fills a Hartogs figure by a holomorphic box
(Proposition~\ref{prop:hartogs}). Ivashkovich's theorem gives a
meromorphic extension to the envelope of holomorphy
\cite{Ivashkovich1992,Ivashkovich2008}. The separate tangencies then
exclude positive-dimensional fibres of its graph. If a first target
coordinate varied along such a fibre, fixing that coordinate together
with $s$ would give an isolated level at a general fibre point. Nearby
ordinary graph points instead lie on one-dimensional levels, because
the same coordinate is independent of $t$. A compact local analytic
level would result, which is impossible in a coordinate ball. The
argument works for every summand rank and for boxes with deficient rank.

The second step crosses a smooth convex boundary at which the complex
tangent has nonzero components in both parameter directions
(Proposition~\ref{prop:mixed}). Bounded projections onto $E_1$ and
$E_2$ compare the normal and tangential speeds of the box. Stokes'
theorem for a K\"ahler form then bounds the areas of shrinking complex
tangent disks. Among many disjoint annuli, one has small area; a local
derivative rescaling argument gives a small image. Product coordinates
on that image build a Hartogs shell whose size exceeds the remaining
distance to the boundary. The first step fills it. Closedness of the
K\"ahler form enters the area estimate, while its metric need not make
the two summands orthogonal.

These extension results combine two holomorphic disk axes into a box
on their entire bidisk. A finite coordinate construction proves
holomorphic dependence on additional parameters, yielding full
polydisk boxes (Section~\ref{sec:boxes}). This is stronger than
continuation of individual germs: after transport, an edge still has
a holomorphic representative on its whole original parameter polydisk.
Consequently a finite rectangular grid combines arbitrary continuous
plaque paths, and a second grid realizes any continuous path as the
diagonal of a rectangle (Section~\ref{sec:global}). The rectangles
continue both the local product map and its inverse. Monodromy and
basepoint-matched covering lifts then give inverse biholomorphisms on
the simply connected covers.

\subsection*{An alternative through collapsed fibres}

The canonical contraction also permits a different continuation method.
We develop it in the appendices because it supplies a useful
neighborhood theorem for singular projective sets as well as another
proof of Corollary~\ref{thm:nef-big}. If $P$ is a compact simply connected
projective set in a smooth projective manifold with complementary integrable tangent
subbundles, and every two points of $P$ are joined by rational chains,
ordinary local crossing extends to a holomorphic map $c:U\times U\to M$
on a neighborhood of the whole set $P$
(Theorem~\ref{thm:crossing}).

The proof first fills products of rational curves, using meromorphic
Hartogs extension and the same complementary-tangency obstruction to
exceptional curves. It then uses proper stable-map spaces
\cite{FP97}, semialgebraic path lifting, and monodromy to continue
\emph{ambient} germs near the possibly singular or reducible set $P$.
For the canonical contraction, Hacon--McKernan's rational-chain theorem
and Takayama's local simple-connectedness theorem provide the required
fibre topology \cite{HM07,Takayama03}. Song's metric completion supplies
a continuous pseudodistance that collapses exactly those fibres.
Crossings are separately nonexpanding for this pseudodistance, and a
finite grid gives path rectangles even for paths contained in a
collapsed fibre. No length bound on those paths is required.

The main proof occupies Sections~\ref{sec:prelim}--\ref{sec:global}.
Section~\ref{sec:canonical} proves the metric estimate and the
nef-and-big application. Appendices~\ref{app:collapsed}--\ref{app:global}
give the alternative continuation mechanism in its dependency order.
All manifolds and varieties are complex. All covers and fundamental
groups are taken in the ordinary topology, and leaves carry their
intrinsic manifold topology.

\section{Product charts and boxes}\label{sec:prelim}

Throughout the proof, $X,E_1,E_2$ satisfy the hypotheses of
Theorem~\ref{thm:main}. Write $\D=\{z\in\C:|z|<1\}$.
Fix a K\"ahler form $\omega$ on $X$ and its Hermitian metric
$g(a,b)=\omega(a,Jb)$ on the real tangent bundle, where $J$ is the
complex structure. We identify the real tangent bundle with the
holomorphic tangent bundle as complex vector bundles, so multiplication
by $i$ on real tangent vectors means application of $J$. The splitting
therefore has complex-linear projections $P_i$ on real tangent vectors.
Compactness gives
\begin{equation}\label{eq:projection-bounds}
\sup_{x\in X}\norm{P_i(x)}<\infty\qquad(i=1,2).
\end{equation}
No orthogonality of the two summands is assumed.

\begin{lemma}\label{lem:product-charts}
Every point of $X$ has a holomorphic coordinate neighborhood with
coordinates $z=(z_1,z_2)$ and range $\D^{r_1}\times\D^{r_2}$, in
which $E_i$ is tangent to the $z_i$-directions. Changes between these
charts are locally of the following form, with each $h_i$ a local
biholomorphism:
\begin{equation}\label{eq:separated-transition}
z'_1=h_1(z_1),\qquad z'_2=h_2(z_2).
\end{equation}
\end{lemma}

\begin{proof}
Holomorphic Frobenius gives a local submersion with kernel $E_i$ for
each $i$. To recall the local construction, choose a coordinate
projection that is nonsingular on the distribution and lift its
coordinate fields to a frame in the distribution. Their brackets
belong to the distribution by integrability and project to zero;
the projection is an isomorphism on the distribution, so the brackets
vanish. The commuting local holomorphic flows applied to a transverse
slice give foliation coordinates and the asserted submersion.

Combine the submersion with kernel $E_2$ and the submersion with kernel
$E_1$. Their combined differential is an isomorphism by complementarity.
The holomorphic inverse function theorem and restriction to a product
of polydisks give the desired coordinates. A chart change preserves
both distributions, so its wrong-block derivatives vanish. On small
product neighborhoods this is exactly
\eqref{eq:separated-transition}, and invertibility of the chart change
makes each block a local biholomorphism.
\end{proof}

We call these \emph{product charts}. A \emph{first plaque} is a slice
with $z_2$ fixed, and a \emph{second plaque} is a slice with $z_1$ fixed.
A continuous path is a first or second \emph{plaque path} if it locally
lies in plaques of the corresponding type. Transitions are required
to be separated only locally; no assertion is made about a single
formula on a disconnected chart overlap.

\begin{definition}\label{def:box}
Let $U\subset\C^a\times\C^b$ be open, with parameter blocks $(s,t)$.
A holomorphic map $f:U\to X$ is a \emph{holomorphic box} if
\[
\dd f(T\C^a\oplus0)\subset E_1,\qquad
\dd f(0\oplus T\C^b)\subset E_2.
\]
The conditions are imposed at every point of $U$. They do not require
$f$ to be an immersion.
\end{definition}

The definition imposes only tangency; the standing integrability
hypothesis supplies the product charts. In such a chart, a box is locally $(f_1(s),f_2(t))$. Conversely,
this local separation implies the two tangencies. The scalar case
$a=b=1$ will be used first. Its parameters $(s,t)$ retain their
prescribed directions even when we describe a domain using other
affine coordinates.

All uses of the identity theorem below are for holomorphic maps or
holomorphic bundle maps on connected domains. In particular, tangency
persists under holomorphic extension: the wrong-summand projections
of the differential vanish wherever the initial box is defined and
therefore throughout the connected extension domain.

\section{Hartogs extension for boxes}\label{sec:hartogs}

Our first extension step is local in the source but uses the compact
K\"ahler target.  A general meromorphic extension can acquire exceptional
fibers.  For a box, the two separate tangencies exclude these fibers, even
when the box has deficient rank.

\begin{proposition}[Hartogs extension]\label{prop:hartogs}
Let $(s,t)$ be the standard coordinates of $\C^2$, and let $(w,v)$ be any
invertible complex-affine coordinate system.  For real numbers
$0<e<a$ and $0<c<b$, put
\[
 H=\bigl\{|w|<e,\ |v|<b\bigr\}
   \ \cup\ \bigl\{|w|<a,\ c<|v|<b\bigr\},
 \qquad
 Q=\bigl\{|w|<a,\ |v|<b\bigr\}.
\]
Every holomorphic map $f:H\to X$ satisfying
\[
 \dd f(\partial_s)\in E_1,
 \qquad
 \dd f(\partial_t)\in E_2
\]
extends uniquely to a holomorphic map $Q\to X$ with the same tangencies.
Thus the affine coordinates defining the Hartogs figure need not be the
coordinates defining the box.
\end{proposition}

\begin{proof}
We first obtain a meromorphic extension.  We then rule out curves in its
graph fibers by comparing two types of level sets: a level isolated at a
putative exceptional curve, and nearby levels containing a parameter
direction.  Finally we show that the remaining finite fibers are singletons.

\emph{Meromorphic extension and tangency.}
The envelope of holomorphy of $H$ is $Q$.  Indeed, expand a scalar
holomorphic function on the shell in a Laurent series in $v$.  The
coefficients of the negative powers are holomorphic functions of $w$;
they vanish for $|w|<e$, and hence for $|w|<a$.  Cauchy estimates on any
intermediate circle $c<|v|<b$ give normal convergence of the resulting
power series on compact subsets of $Q$.  This is the usual Hartogs
extension calculation; since the bidisk $Q$ is a domain of holomorphy, it
also identifies the envelope.

Ivashkovich's extension theorem states that a meromorphic map from a
domain in a Stein manifold to a compact K\"ahler manifold extends
meromorphically to the envelope of holomorphy
\cite[Theorem~2.2]{Ivashkovich2008}; see also \cite{Ivashkovich1992}.
Here $H$ is a domain in the Stein manifold $\C^2$, and $f$ is holomorphic,
so the theorem gives a meromorphic extension to $Q$.

Let $\Gamma\subset Q\times X$ be its graph, with its reduced analytic
structure, and let $\rho:\Gamma\to Q$ be the first projection.  The graph
is an irreducible analytic surface, $\rho$ is proper, and there is a proper
analytic subset $A\subset Q$, disjoint from $H$, such that over
$U=Q\setminus A$ the graph is the ordinary graph of a holomorphic map
$F:U\to X$.  This ordinary graph is dense in $\Gamma$.  The set $U$ is
connected: the complement of a proper complex analytic subset of a
connected complex manifold is connected.  Consequently, the identity
theorem applied to the two holomorphic bundle-valued expressions gives
\begin{equation}\label{eq:meromorphic-box-tangency}
 P_2\dd F(\partial_s)=0,
 \qquad
 P_1\dd F(\partial_t)=0
 \quad\hbox{on }U.
\end{equation}
Here $P_i:T_X\to E_i$ is the holomorphic projection determined by the
specified splitting.

\emph{No positive-dimensional fiber.}
We use the local irreducible decomposition, dimension theorem, and
regular-locus properties of complex analytic sets; see
\cite[Chapter~II, Section~4]{DemaillyCADG}.
We record first a point that is important when $\Gamma$ is singular or
locally reducible.  On every local irreducible surface branch of $\Gamma$,
both source coordinates $s$ and $t$ are nonconstant.  To see this, choose
a relatively open part of that branch away from the other local branches.
It meets the dense ordinary graph.  Near such a point, $\rho$ is an
isomorphism onto an open subset of $Q$, and $s,t$ are source coordinates.
Neither coordinate can therefore be constant on the branch.  It follows
that, for every $s_0$, the analytic set
\begin{equation}\label{eq:graph-coordinate-level}
 \Gamma\cap\{s=s_0\}
\end{equation}
has local dimension at most one; the same holds with $t$ in place of $s$.

Suppose a fiber over $(s_0,t_0)$ has positive dimension.  Since it is
contained in \eqref{eq:graph-coordinate-level}, it has a curve component
$C$.  Take a small piece of $C$ in a product chart
$z=(z_1,z_2)$ of $X$.  Some scalar component $h$ of $z_1$ or $z_2$ is
nonconstant on this piece: otherwise the fiber curve, which lies in
$\{(s_0,t_0)\}\times X$, would be a point.

First suppose $h$ belongs to $z_1$.  The curve piece is a component of
the reduced set \eqref{eq:graph-coordinate-level}, since that set has
dimension at most one.  Choose a general point $q$ of the
curve piece such that the reduced set \eqref{eq:graph-coordinate-level}
is a smooth curve near $q$ and $h$ is nonconstant on its germ.  This is
possible because the one-dimensional analytic set has locally finitely
many curve components, distinct components meet discretely, and the
singular points of a reduced curve are discrete.  In particular, $q$ is
an isolated point of the level set
\[
 (s,h)^{-1}\bigl(s_0,h(q)\bigr)\subset\Gamma.
\]
This assertion concerns the reduced coordinate level, and does not
require $\Gamma$ itself to be smooth along $C$.

Work in ambient complex coordinates on $Q\times X$ around $q$, and choose
a sufficiently small closed Euclidean ball $\overline B$ centered at $q$.
Its closure lies in this coordinate neighborhood, and the indicated
level set misses $\Gamma\cap\partial B$.  The image of this boundary
intersection under the holomorphic map
\[
 \Psi=(s,h)
\]
is compact and does not contain $\Psi(q)$.  Choose points $q_j$ of the
ordinary graph tending to $q$.  For all sufficiently large $j$, the level
value $\Psi(q_j)$ also misses that compact image.  Therefore
\[
 K_j=\Gamma\cap\overline B\cap\Psi^{-1}\bigl(\Psi(q_j)\bigr)
\]
is a compact analytic subset of the open ball $B$: it has no points on
$\partial B$, so there is no boundary truncation of the analytic level.
On the ordinary graph near $q_j$, however,
\eqref{eq:meromorphic-box-tangency} says that $h$ is independent of $t$.
Varying $t$ with $s$ fixed gives a one-dimensional germ in $K_j$ through
$q_j$.  This contradicts the fact that a compact analytic subset of an
open set in Euclidean complex space has no positive-dimensional
component.  One way to see the latter fact is to maximize the strictly
plurisubharmonic squared coordinate norm on such a component and restrict
it to a nonconstant analytic curve germ through a maximum point.

If $h$ belongs to $z_2$, the identical argument uses $\Psi=(t,h)$ and
varies $s$.  Thus every fiber of $\rho$ has dimension zero.  Properness
then makes each fiber finite.

\emph{Singleton fibers and holomorphic removal.}
Fix $p\in Q$, and write its fiber as
$\{(p,x_1),\ldots,(p,x_k)\}$.  Choose pairwise disjoint target
neighborhoods $V_1,\ldots,V_k$ with pairwise disjoint closures and
$x_i\in V_i$.  Properness implies that, after shrinking a base ball
$D$ about $p$, every point of $\rho^{-1}(D)$ has its target coordinate in
$\bigcup_iV_i$.  Otherwise a sequence of graph points over points tending
to $p$, outside this union, would have a limit in the central fiber.
The ordinary locus $D\setminus A$ is connected, so its image under $F$
lies in a single $V_i$.  Density of the ordinary graph then places every
$x_j$ in $\overline V_i$.  The disjoint closures imply $k=1$.

The proper bijection $\rho$ is a homeomorphism.  Its inverse followed by
the second projection defines a continuous extension
$\widehat F:Q\to X$.  Near each source point its image lies in a target
coordinate chart.  The coordinate functions are continuous and
holomorphic off $A$, hence holomorphic by the removable-singularity
theorem for locally bounded holomorphic functions across analytic
subsets of a complex manifold.  Thus $\widehat F$ is holomorphic.
Equation~\eqref{eq:meromorphic-box-tangency} extends by the identity
theorem, giving the required box tangencies.  Uniqueness follows from
the identity theorem on the connected bidisk $Q$.
\end{proof}

\section{Extension across a mixed boundary}\label{sec:boundary}

We next extend a box across a smooth boundary point whose complex tangent
involves both source directions.  The K\"ahler form controls the areas of
shrinking slices.  An annulus of small area then supplies the shell of a
Hartogs figure large enough to cross the boundary.  We first isolate the
elementary compactness argument that converts small area into a small
image. Its proof is a local form of the derivative normalization used in
Brody's reparametrization argument \cite{Brody78}, with vanishing area
providing the contradiction.

\begin{lemma}[Small area]\label{lem:small-area}
Let $M$ be a compact Hermitian manifold, with associated positive
$(1,1)$-form $\omega_M$, and let $U\subset\C$ be a domain.  Suppose that
$h_j:U\to M$ are holomorphic and
\[
  \int_U h_j^*\omega_M\longrightarrow 0.
\]
Then $\norm{\dd h_j}$ tends to zero uniformly on every compact subset of $U$.
Here the source has its Euclidean metric.  Closedness of $\omega_M$ is not
required.
\end{lemma}

\begin{proof}
It suffices to work on a disk $D$ with $\overline D\subset U$.  Put
$d(z)=\operatorname{dist}(z,\partial D)$ and
\[
  M_j=\max_{\overline D}d(z)\norm{\dd h_j(z)}.
\]
If $M_j\not\to0$, pass to a subsequence with $M_j\ge\varepsilon>0$.
Choose a maximizing point $z_j\in D$, and set
$a_j=\norm{\dd h_j(z_j)}$ and $d_j=d(z_j)$, so $a_jd_j=M_j$.
The maps
\[
  k_j(\zeta)=h_j(z_j+\zeta/a_j)
\]
are defined for $|\zeta|<\varepsilon/4$.  On this disk,
$d(z_j+\zeta/a_j)\ge d_j/2$, and maximality gives
\[
  \norm{\dd k_j(\zeta)}\le2,
  \qquad \norm{\dd k_j(0)}=1.
\]
After taking a subsequence, $k_j(0)$ converges to a point of $M$.
The derivative bound puts the images of a fixed smaller disk in a
relatively compact coordinate neighborhood of that point.  Montel's
theorem and the Cauchy estimates give a further subsequence converging
there together with its derivatives.  Its limit has derivative norm $1$
at the origin and therefore positive area on a smaller disk.  This
contradicts the vanishing areas of $k_j$, which are bounded by the areas
of $h_j$ and are unchanged by the affine reparametrization.  Thus
$M_j\to0$.  On each smaller disk $d$ is bounded below, so the derivatives
tend uniformly to zero.  A finite covering proves the assertion on any
compact subset of $U$.
\end{proof}

\begin{proposition}[Mixed-boundary extension]\label{prop:mixed}
Let $B\subset\C^2_{s,t}$ be a convex open set, and let $f:B\to X$ be a
holomorphic box.  Suppose that $p\in\partial B$ has a $C^2$ defining
function with nonzero differential.  If the complex tangent line to
$\partial B$ at $p$ is generated by a vector with both components
nonzero, then $f$ extends as a holomorphic box to a neighborhood of $p$,
agreeing with $f$ on its intersection with $B$.
\end{proposition}

\begin{proof}
Use the fixed K\"ahler form $\omega$ and its Hermitian norm, and recall
the bounded projections $P_i$ from \eqref{eq:projection-bounds}.

Write $p=(s_*,t_*)$.  Choose complex-affine coordinates $(u,v)$ with
\begin{equation}\label{eq:boundary-affine}
  (s,t)=(s_*,t_*)+(a_1,a_2)u+(b_1,b_2)v,
  \qquad b_1b_2\ne0,
\end{equation}
where the $v$-direction is complex tangent and the positive real
$u$-direction points outward.  All constants below may depend on this
boundary point and these coordinates.

\paragraph{Shrinking slices and bounded area.}
Take a local defining function $\rho$ with $B=\{\rho<0\}$ near $p$.
For real $x$ and small $v$, Taylor's theorem gives
\[
  \rho(x,v)\le c x+M(x^2+|v|^2),\qquad c>0,
\]
because both real derivatives in the complex tangent direction vanish.
Choose $k>0$ sufficiently small.  For all $x<0$ sufficiently close to
zero, the closed disks
\[
  D_x=\{u=x,\ |v|\le r(x)\},\qquad r(x)=k\sqrt{-x},
\]
are compactly contained in $B$.

Let $e_u=\partial_{\operatorname{Re}u}$ and
$e_v=\partial_{\operatorname{Re}v}$.  The box tangencies and complex
linearity give the exact identity
\[
  \dd f(e_u)=\sum_{i=1}^2\frac{a_i}{b_i}P_i\dd f(e_v).
\]
Consequently
\begin{equation}\label{eq:boundary-speed}
  \norm{\dd f(e_u)}\le C\norm{\dd f(e_v)},\qquad
  C=\sum_{i=1}^2\left|\frac{a_i}{b_i}\right|
       \sup_X\norm{P_i}.
\end{equation}
Holomorphicity makes $\norm{\dd f(e_v)}$ equal to the speed in every unit
real direction of the $v$-plane, including the radial direction.

Set $\Omega=f^*\omega$ and $A(x)=\int_{D_x}\Omega$.  Closedness of
$\Omega$ and Stokes' theorem, applied to the moving disks, yield
\[
 A'(x)=\int_0^{2\pi}
   \Omega\bigl(e_u+r'(x)\partial_r,\partial_\theta\bigr)
       \big|_{u=x,\ v=r(x)e^{i\theta}}\,\dd\theta.
\]
Indeed the boundary velocity is $e_u+r'(x)\partial_r$, and the circles
have their counterclockwise orientation.  On a boundary circle put
$q=\norm{\dd f(e_v)}$.  Since $\partial_\theta=rJ\partial_r$,
\[
  \Omega(\partial_r,\partial_\theta)=rq^2,
  \qquad
  |\Omega(e_u,\partial_\theta)|\le Crq^2.
\]
Thus
\[
  A'(x)\le\int_0^{2\pi}r(x)\bigl(C+r'(x)\bigr)q^2\,\dd\theta.
\]
Here $r'(x)=-k/(2\sqrt{-x})$.  Choose $x_0<0$ sufficiently close to
zero that $r'(x)\le-C$ for $x_0\le x<0$.  We obtain the uniform bound
\begin{equation}\label{eq:boundary-area}
  0\le A(x)\le A(x_0)<\infty\qquad(x_0\le x<0).
\end{equation}
No product property of the K\"ahler metric has been used.

\paragraph{A small-image annulus.}
Let $x_j=-\delta_j$ with $\delta_j\downarrow0$.  Among the scales
$\lambda=5^\ell\delta_j^{3/4}$, with integer $\ell\ge0$, in
$[\delta_j^{3/4},\delta_j^{5/8}]$, the annuli
\[
  \lambda/2<|v|<2\lambda
\]
are pairwise disjoint, and their number tends to infinity.  For large
$j$ they lie in $D_{x_j}$.  By \eqref{eq:boundary-area}, one of them,
with scale $\lambda_j$, has area tending to zero.  These choices satisfy
\begin{equation}\label{eq:boundary-scales}
  \frac{\delta_j}{\lambda_j}\le\delta_j^{1/4}\longrightarrow0,
  \qquad
  \frac{\lambda_j}{\sqrt{\delta_j}}
       \le\delta_j^{1/8}\longrightarrow0.
\end{equation}
Apply Lemma~\ref{lem:small-area} to
$h_j(\zeta)=f(x_j,\lambda_j\zeta)$ on $1/2<|\zeta|<2$.
The derivatives tend uniformly to zero on a neighborhood of
$3/4\le|\zeta|\le3/2$.  This compact annulus has bounded intrinsic
diameter, so the image of
\[
  u=x_j,\qquad 3\lambda_j/4<|v|<3\lambda_j/2
\]
has diameter tending to zero.  A finite cover of $X$ by product charts
has a Lebesgue number.  Hence, for large $j$, this entire image lies in
one product chart $z:V\to V_1\times V_2$.  Write its coordinate functions
as $(g_1(v),g_2(v))$.

Figure~\ref{fig:boundary-scales} separates the two scale comparisons.
\begin{figure}[tbp]
\centering
\setlength{\unitlength}{1pt}
\begin{picture}(430,145)
  \put(100,125){\makebox(0,0){\small normal real slice ($v=0$)}}
  \put(15,75){\vector(1,0){180}}
  {\color{blue!65!black}\linethickness{2pt}\put(35,75){\line(1,0){130}}}\ignorespaces
  \put(100,71){\line(0,1){8}}
  \put(116,68){\line(0,1){14}}
  \put(100,75){\circle*{3}}
  \put(116,75){\circle*{3}}
  \put(91,58){\makebox(0,0){\small $x_j=-\delta_j$}}
  \put(127,91){\makebox(0,0){\small $p: u=0$}}
  \put(35,101){\line(1,0){130}}
  \put(35,98){\line(0,1){6}}
  \put(165,98){\line(0,1){6}}
  \put(100,111){\makebox(0,0){\small shell width $2\eta\lambda_j$}}
  \put(193,63){\makebox(0,0){\small $u$}}
  \put(108,37){\makebox(0,0){\small $\delta_j/\lambda_j\longrightarrow0$}}
  \put(320,125){\makebox(0,0){\small complex tangent slice ($u=x_j$)}}
  \put(320,69){\circle{90}}
  {\color{blue!65!black}\put(320,69){\circle{34}}\put(320,69){\circle{43}}}\ignorespaces
  \put(320,69){\circle*{2}}
  \put(320,69){\vector(1,0){45}}
  \put(376,58){\makebox(0,0)[l]{\small $r(x_j)$}}
  \put(320,69){\vector(-4,3){16}}
  \put(300,90){\makebox(0,0)[r]{\small $\lambda_j$}}
  \put(320,8){\makebox(0,0){\small $\lambda_j/\sqrt{\delta_j}\longrightarrow0$}}
\end{picture}
\caption{The two scale inequalities used to cross the boundary. On the
left, the filled Hartogs domain extends in the $u$ direction farther than
the distance $\delta_j$ from the slice to $p$. On the right, its shell
has radius comparable to $\lambda_j$ and fits well inside the slice disk
of radius $r(x_j)=k\sqrt{\delta_j}$. The annulus has a small image in one
target product chart. These are separate source slices in the affine
coordinates $(u,v)$, drawn schematically and not at a common scale.}
\label{fig:boundary-scales}
\end{figure}


We have found an annulus whose radius is much larger than the remaining
distance $\delta_j$ to the boundary.  We now extend it transversely to
make a Hartogs shell.

\paragraph{The shell and its agreement with the old box.}
Put $w=u-x_j$ and $c_i=a_i/b_i$.  Choose $\eta>0$, independent of $j$,
so that $\eta\max(1,|c_1|,|c_2|)<1/8$.  On
\[
  S_j=\{|w|<\eta\lambda_j,\quad
                \lambda_j<|v|<5\lambda_j/4\}
\]
define
\begin{equation}\label{eq:boundary-shell}
  F_j(w,v)=z^{-1}\bigl(g_1(v+c_1w),g_2(v+c_2w)\bigr).
\end{equation}
Both shifted arguments stay in
$3\lambda_j/4<|v|<3\lambda_j/2$.  The range of $z$ is a product, so
combining the two valid factor coordinates in
\eqref{eq:boundary-shell} needs no further target-chart margin.  Moreover,
\[
  v+c_1w=\frac{s-s_*-a_1x_j}{b_1},\qquad
  v+c_2w=\frac{t-t_*-a_2x_j}{b_2}.
\]
Thus $F_j$ is a box in the original $(s,t)$ directions.

The two maps agree on an open germ at every point of $S_j$ with $w=0$.
To see this, continuity puts the old map $f$ in $V$ on a neighborhood
of that point.  There its separated expression is
$(\varphi_1(s),\varphi_2(t))$.  Restriction to $u=x_j$ gives the local
expressions for $g_i$, and the displayed identities show that
\eqref{eq:boundary-shell} reproduces $f$ on a full neighborhood.  This
is stronger than equality merely on the slice $w=0$.

The closed slice disk $u=x_j$, $|v|\le5\lambda_j/4$ is contained in
$B$.  Hence there is $0<e_j<\eta\lambda_j$ such that the original map
is defined on
\[
  C_j=\{|w|<e_j,\ |v|<5\lambda_j/4\}.
\]
The overlap $C_j\cap S_j$ is a connected product of a disk and an
annulus.  The identity theorem for maps into $X$ therefore extends the
open-germ agreement to this entire overlap, even if the old map elsewhere
leaves $V$.  We have a box on the Hartogs figure $C_j\cup S_j$.

By Proposition~\ref{prop:hartogs}, it extends to a box on
\[
  Q_j=\{|w|<\eta\lambda_j,\ |v|<5\lambda_j/4\}.
\]
Both $Q_j$ and $B$ are convex in the underlying real affine space.
Their intersection is connected and contains $C_j$, so the extension
agrees with $f$ throughout $Q_j\cap B$.  Finally $p$ has coordinates
$(w,v)=(\delta_j,0)$, which belong to $Q_j$ for large $j$ by
\eqref{eq:boundary-scales}.  This gives the required extension across
$p$.
\end{proof}

\section{Boxes with prescribed axes}\label{sec:boxes}

The mixed-boundary extension result lets us combine two holomorphic maps
on full disks, not just their germs. We then establish the parameter
dependence needed to combine maps from higher-dimensional polydisks.

\begin{proposition}[Disk axes]\label{prop:disk-axes}
Let \(\alpha,\beta:\D\to X\) be holomorphic maps with
\(\dd\alpha(T_\D)\subset E_1\), \(\dd\beta(T_\D)\subset E_2\), and
\(\alpha(0)=\beta(0)\). There is a unique holomorphic box
\(F:\D^2\to X\) such that
\[
 F(s,0)=\alpha(s),\qquad F(0,t)=\beta(t).
\]
\end{proposition}

\begin{proof}
Product coordinates at the common initial point give a unique box germ:
take the first coordinate from \(\alpha\) and the second from \(\beta\).
Any two boxes extending this germ agree on a connected common domain by
the identity theorem.

Fix an integer \(m\geq1\), and consider the convex domains
\[
 B_m(R)=\{(s,t):|s|^{2m}+|t|^{2m}<R^{2m}\},\qquad 0<R\leq1.
\]
Let \(R_*\) be the supremum of the radii to which the germ extends as
a box. The nested extensions agree and give a box on \(B_m(R_*)\).
Its values on the two axes are the prescribed maps, by the one-variable
identity theorem. Suppose that \(R_*<1\).

At a boundary point with \(st\neq0\), the complex tangent line has both
components nonzero: both coefficients of
\(\partial(|s|^{2m}+|t|^{2m})\) are nonzero there. Proposition~\ref{prop:mixed}
therefore extends the box across that point.
At a boundary point \((s_*,0)\), choose a disk \(U\) around \(s_*\),
with \(\overline U\subset\D\), such that \(\alpha(U)\) lies in a
product chart \(z=(z_1,z_2)\). Choose \(s'\in U\) with
\(|s'|<R_*\). For sufficiently small \(\epsilon>0\), the formula
\[
 G(s,t)=z^{-1}\bigl(z_1(\alpha(s)),z_2(F(s',t))\bigr),
 \qquad (s,t)\in U\times\D_\epsilon,
\]
is defined and is a box. Here \(\D_\epsilon=\{|t|<\epsilon\}\).
The product range of the chart makes the formula well-defined. Local
separation shows that \(G=F\) near \((s',0)\). Their common domain
\((U\times\D_\epsilon)\cap B_m(R_*)\) is convex, so they agree
throughout it. The other axis is treated symmetrically.

These boundary extensions glue. Indeed, restrict each one to an ordinary
Euclidean ball centered at its boundary point. If two such balls intersect,
their intersection contains a point of the segment joining their centers.
That segment lies in \(\overline{B_m(R_*)}\) by convexity. Since the
intersection is open, it consequently meets \(B_m(R_*)\). The two
extensions agree there, hence on their connected intersection. A finite
subcollection covers the compact boundary, giving an extension to a
neighborhood of \(\overline{B_m(R_*)}\). This neighborhood contains
\(B_m(R')\) for some \(R_*<R'<1\), a contradiction.

Thus the germ extends to \(B_m(1)\) for every \(m\). These domains
increase to \(\D^2\), and the resulting boxes agree on their connected
overlaps. Their union is the required box. The initial uniqueness
observation proves uniqueness on \(\D^2\).
\end{proof}

We will also use boxes with real parameters. A \emph{continuous box} on
a product of closed intervals is a continuous map into \(X\) that is
locally separated in product coordinates. Equivalently, every horizontal
slice locally follows an \(E_1\)-plaque and every vertical slice locally
follows an \(E_2\)-plaque. If a whole subrectangle maps into one product
chart, its first coordinate depends only on the horizontal parameter
and its second only on the vertical parameter: the unwanted coordinate
is locally constant along each connected slice.

\begin{lemma}[Uniqueness of continuous boxes]\label{lem:box-unique}
Two continuous boxes on \([0,1]^2\) with the same bottom and left edges
are equal.
\end{lemma}

\begin{proof}
Choose a Lebesgue number \(\lambda>0\) for a finite product-chart
cover of the compact metric space \(X\). Uniform continuity gives one
finite rectangular grid such that the image of each closed tile under
either map has diameter less than \(\lambda/3\). Whenever the two maps
agree at a tile's lower-left corner, the union of their tile images has
diameter less than \(2\lambda/3\), so both images lie in one product
chart. In that chart each box is determined by its bottom first
coordinate and left second coordinate. Starting with the prescribed
edges and proceeding by increasing sum of the tile indices proves
equality on the whole grid.
\end{proof}

\begin{proposition}[Polydisk axes]\label{prop:polydisk-axes}
Let \(S\) and \(T\) be unit polydisks of positive dimensions. Suppose
\(\alpha:S\to X\) and \(\beta:T\to X\) are holomorphic maps tangent
to \(E_1\) and \(E_2\), respectively, and
\(\alpha(p_0)=\beta(q_0)\) for some \((p_0,q_0)\in S\times T\).
There is a unique holomorphic box \(F:S\times T\to X\) with
\[
 F(p,q_0)=\alpha(p),\qquad F(p_0,q)=\beta(q).
\]
\end{proposition}

\begin{proof}
Polydisk automorphisms reduce the problem to \(p_0=q_0=0\).
For fixed \((p,q)\), the maps
\[
 \sigma\longmapsto\alpha(\sigma p),\qquad
 \tau\longmapsto\beta(\tau q)
\]
are holomorphic on disks of radii greater than one. After rescaling,
Proposition~\ref{prop:disk-axes} gives a box whose restriction to the
closed real square is denoted by \(R_{p,q}\). Set
\[
 F(p,q)=R_{p,q}(1,1).
\]
The constant-axis cases and Lemma~\ref{lem:box-unique} give the prescribed
axis values. It remains to prove holomorphicity and the two tangencies;
pointwise existence of \(R_{p,q}\) alone does not give parameter
continuity.

Fix \((p_*,q_*)\), and choose a finite grid whose closed tile images
under \(R_{p_*,q_*}\) lie in product charts. We reconstruct this square
for nearby \((p,q)\), using these fixed charts. Begin with its two
given axes, and process tiles from the bottom and left. The induction
maintains the following properties: each constructed tile depends
jointly continuously on its real parameters and on \((p,q)\); it is a
continuous box in the real parameters; and, at each fixed pair of real
parameters, it is holomorphic in \((p,q)\), with the \(p\)-directions
in \(E_1\) and the \(q\)-directions in \(E_2\). These properties hold
on the initial axes.

To fill one tile, denote its already constructed bottom and left edges
by \(B(\sigma;p,q)\) and \(L(\tau;p,q)\). At the reference parameter
their images lie in the selected product chart \(z=(z_1,z_2)\).
Joint continuity and compactness of those edges put both images in that
same chart after shrinking the parameter neighborhood. Define the tile
by
\begin{equation}\label{eq:parameter-tile}
 z^{-1}\bigl(z_1(B(\sigma;p,q)),z_2(L(\tau;p,q))\bigr).
\end{equation}
The edges meet at a common corner. Plaque constancy on each edge shows
that this formula reproduces both of them. It is jointly continuous and
is separated in \(\sigma,\tau\). For each fixed \(\sigma,\tau\), it
is holomorphic in \((p,q)\); moreover, the induction hypothesis gives
\[
 \partial_q\bigl(z_1\circ B\bigr)=0,
 \qquad
 \partial_p\bigl(z_2\circ L\bigr)=0.
\]
Thus the two parameter tangencies are preserved as well. At
\((p_*,q_*)\), the reconstructed tile is the original tile, so the
induction continues. There are only finitely many tiles and therefore
only finitely many shrinkings.

The tile maps agree on common edges and form a continuous box across
the seams. To check the latter assertion at a seam or a vertex, use a
product chart around its common image. On a sufficiently small real
neighborhood, the unwanted coordinate along a horizontal or vertical
slice is constant on each tile piece and agrees at their endpoints;
it is therefore constant across the seam. By
Lemma~\ref{lem:box-unique}, the reconstructed square is \(R_{p,q}\)
for every nearby parameter pair. Its upper-right value is consequently
\(F(p,q)\), proving that \(F\) is a holomorphic box near
\((p_*,q_*)\), and hence everywhere. Finally, the axes determine its
germ at \((0,0)\), so the identity theorem proves uniqueness.
\end{proof}

\section{Path rectangles and the universal cover}\label{sec:global}

We now turn the holomorphic box theorem into continuation along arbitrary
continuous paths. The important point is that a transported edge retains
a holomorphic parametrization on an entire polydisk. Its image need not
remain small or lie in a single original chart.

\begin{proposition}[Path rectangles]\label{prop:rectangles}
The following assertions hold.
\begin{enumerate}
\item If \(\alpha,\beta:[0,1]\to X\) locally follow \(E_1\)- and
\(E_2\)-plaques, respectively, and \(\alpha(0)=\beta(0)\), there is
a continuous box \(R:[0,1]^2\to X\) with
\(R(s,0)=\alpha(s)\) and \(R(0,t)=\beta(t)\).
\item Every continuous path \(\gamma:[0,1]\to X\) is the diagonal
of a continuous box: \(R(s,s)=\gamma(s)\).
\end{enumerate}
\end{proposition}

\begin{proof}
For the first assertion, subdivide the horizontal interval into pieces
\(I_i=[u_{i-1},u_i]\) and the vertical interval into pieces
\(J_j=[v_{j-1},v_j]\) on which the respective paths lie in single
plaques. Each horizontal piece has a representation
\[
 \alpha(s)=A_i(a_i(s)),\qquad s\in I_i,
\]
where \(a_i:I_i\to\D^{r_1}\) is continuous and
\(A_i:\D^{r_1}\to X\) is a holomorphic plaque parametrization
tangent to \(E_1\). Similarly write
\(\beta(t)=B_j(b_j(t))\), with
\(b_j:J_j\to\D^{r_2}\) continuous and \(B_j\) holomorphic and
tangent to \(E_2\). Here \(r_i=\rank E_i\).

Fill the tiles \(I_i\times J_j\) from the bottom and left. Inductively,
the known edges of a tile have the forms
\(A(a_i(s))\) and \(B(b_j(t))\), where \(A\) and \(B\) are
holomorphic on the \emph{entire} polydisks \(\D^{r_1}\) and
\(\D^{r_2}\), tangent to the respective bundles. They agree at their
common corner. Apply Proposition~\ref{prop:polydisk-axes} with origins
\(a_i(u_{i-1})\) and \(b_j(v_{j-1})\), obtaining a holomorphic box
\(G:\D^{r_1}\times\D^{r_2}\to X\). Its evaluation
\[
 R(s,t)=G(a_i(s),b_j(t))
\]
fills the tile. The opposite horizontal and vertical edges have the
holomorphic representatives
\[
 a\longmapsto G(a,b_j(v_j)),\qquad
 b\longmapsto G(a_i(u_i),b),
\]
again on the whole original polydisks. This proves the induction without
any shrinking of the edge domains.

For the second assertion, choose a subdivision
\(0=u_0<\cdots<u_N=1\) such that each
\(\gamma(I_i)\), \(I_i=[u_{i-1},u_i]\), lies in a product chart
with inverse \(\theta_i:\D^{r_1}\times\D^{r_2}\to X\).
Write \(\gamma(s)=\theta_i(a_i(s),b_i(s))\) on that interval.
Define the diagonal tile by
\[
 R(s,t)=\theta_i(a_i(s),b_i(t)),\qquad (s,t)\in I_i^2.
\]
These tiles agree at consecutive diagonal vertices, and their edges have
the full-polydisk representations used above.

Fill the remaining tiles \(I_i\times I_j\) in successive bands away
from the diagonal, that is, in increasing order of \(|i-j|\).
If \(i<j\), use the already filled bottom and right
edges; if \(i>j\), use the top and left edges. The same polydisk box
construction works from any corner, since the origins in
Proposition~\ref{prop:polydisk-axes} are arbitrary. The two input edges
have matching corner values: for \(|i-j|=1\) this is the prescribed
diagonal value, and for \(|i-j|>1\) it follows from the intervening
tile closer to the diagonal. More explicitly, when \(i<j-1\), the
bottom-right corner is shared with the previously filled tile
\(I_{i+1}\times I_{j-1}\). The case \(i>j+1\) is symmetric.
Tiles filled at the same stage share no edge, and all their shared
vertices already lie on prescribed input edges. Thus every tile is
consistent with the previously assigned values, and its output edges
retain their full-polydisk representatives.

In either construction, the finite collection of tiles glues
continuously. The local plaque condition holds across seams for the
same reason as in the proof of Proposition~\ref{prop:polydisk-axes}:
in a product chart near a seam point, the unwanted coordinate is
constant on consecutive slice pieces and agrees at their common
endpoints. The resulting map is a continuous box.
\end{proof}

Figure~\ref{fig:grid} illustrates both filling orders. The transported
edges retain their full-polydisk representatives; this, rather than a
bound on their image sizes, is what allows each new tile to be filled.
\begin{figure}[tbp]
\centering
\setlength{\unitlength}{1pt}
\begin{picture}(430,175)
  \put(102,163){\makebox(0,0){\small (a) prescribed axes}}
  \put(322,163){\makebox(0,0){\small (b) prescribed diagonal}}
  {\color{black!12}
   \put(55,45){\rule{30pt}{30pt}}
   \put(275,45){\rule{30pt}{30pt}}
   \put(305,75){\rule{30pt}{30pt}}
   \put(335,105){\rule{30pt}{30pt}}}\ignorespaces
  {\color{blue!12}
   \put(85,45){\rule{30pt}{30pt}}
   \put(305,45){\rule{30pt}{30pt}}}\ignorespaces
  \thinlines
  \multiput(55,45)(30,0){4}{\line(0,1){90}}
  \multiput(55,45)(0,30){4}{\line(1,0){90}}
  \multiput(275,45)(30,0){4}{\line(0,1){90}}
  \multiput(275,45)(0,30){4}{\line(1,0){90}}
  {\color{black!55}\put(275,45){\line(1,1){90}}}\ignorespaces
  \thicklines
  {\color{blue!65!black}
   \put(55,45){\line(1,0){90}}
   \put(85,45){\vector(1,0){30}}
   \put(305,75){\vector(1,0){30}}}\ignorespaces
  {\color{red!65!black}
   \put(55,45){\line(0,1){90}}
   \put(85,45){\vector(0,1){30}}
   \put(305,75){\vector(0,-1){30}}}\ignorespaces
  \put(85,45){\circle*{3}}
  \put(305,75){\circle*{3}}
  \put(100,60){\makebox(0,0){\small $R$}}
  \put(320,60){\makebox(0,0){\small $R$}}
  \put(100,30){\makebox(0,0){\small $\alpha(s)$}}
  \put(47,90){\makebox(0,0)[r]{\small $\beta(t)$}}
  \put(155,45){\makebox(0,0){\small $s$}}
  \put(55,144){\makebox(0,0){\small $t$}}
  \put(375,45){\makebox(0,0){\small $s$}}
  \put(275,144){\makebox(0,0){\small $t$}}
  \put(320,29){\makebox(0,0){\small $R(s,s)=\gamma(s)$}}
  {\color{black!12}\put(105,5){\rule{9pt}{9pt}}}\ignorespaces
  \put(120,9){\makebox(0,0)[l]{\small completed tiles}}
  {\color{blue!12}\put(255,5){\rule{9pt}{9pt}}}\ignorespaces
  \put(270,9){\makebox(0,0)[l]{\small next tile}}
\end{picture}
\caption{The two grid constructions in the parameter square. In (a),
the bottom and left axes are prescribed, and tiles are filled in
increasing order of the sum of their indices. In (b), the diagonal tiles
are prescribed first, and the remaining tiles are filled in successive
bands away from the diagonal. In each highlighted tile, the blue
horizontal edge and red vertical edge are already known and meet at
the marked corner. The box with these axes supplies the other two edges.
The reversed red arrow in (b) indicates that the same construction can
start at any corner. No target-space size estimate is depicted.}
\label{fig:grid}
\end{figure}


Fix a basepoint \(o\in X\). For \(i=1,2\), let \(L_i\) be the
leaf through \(o\), endowed with its intrinsic plaque topology, not
the subspace topology inherited from \(X\).

\begin{lemma}[Intrinsic leaves]\label{lem:leaves}
Each \(L_i\) is a connected, Hausdorff, second-countable complex
manifold of dimension \(r_i\), and its inclusion into \(X\) is an
injective holomorphic immersion. It admits a connected simply connected
cover \(\widetilde L_i\), also a second-countable complex manifold.
A path locally following \(E_i\)-plaques is continuous in the intrinsic
topology of its leaf.
\end{lemma}

\begin{proof}
The leaf consists of points reachable from \(o\) by finite chains of
\(E_i\)-plaques. Plaque coordinates give its complex atlas: near a
point of intersection, two plaques through that point coincide, and
their coordinate change is holomorphic. The inclusion into \(X\) is
an injective continuous immersion. In particular, distinct points of
the leaf are separated by inverse images of disjoint ambient open sets,
so the leaf is Hausdorff. The definition by plaque chains gives path
connectedness and the asserted continuity of plaque paths.

For second countability, take a countable ambient product atlas. Start
with a plaque through \(o\). For each plaque already reached and each
ambient chart, their intersection has at most countably many connected
components: it is an open subset of a second-countable plaque. Each
component lies in a single plaque of the new chart, since the transverse
coordinate is locally constant. Adjoin these plaques and repeat.
There are only countably many plaques after all finite stages, and
every plaque path from \(o\) is covered by a finite chain of this
kind. Their countable coordinate bases form a countable basis for the
leaf.

The usual covering-space construction applies because a manifold is
locally path connected and semilocally simply connected. Its universal
cover is second countable as well. One way to see the latter detail is
to refine the countable atlas to simply connected coordinate disks.
There are countably many components of their pairwise overlaps. A
based loop is described, up to homotopy, by a finite sequence of these
charts and overlap components; fixing paths within the charts and
components supplies representatives for all such descriptions. Hence
the fundamental group is countable. The lifted countable chart basis
has only countably many sheets and is therefore a countable basis for
the universal cover. Lifting the holomorphic charts gives its complex
structure.
\end{proof}

A product chart at \(o\) supplies mutually inverse germs
\[
 H_o:(L_1\times L_2,(o,o))\longrightarrow(X,o),
 \qquad
 K_o:(X,o)\longrightarrow(L_1\times L_2,(o,o)).
\]
The first combines coordinates on the two plaques through \(o\); the
second projects onto them. Separated chart transitions show that the
germs do not depend on the chosen product chart. The next result gives
analytic continuation in both directions.

\begin{proposition}[Transport of inverse germs]\label{prop:germ-transport}
Let \(R:[0,1]^2\to X\) be a continuous box with \(R(0,0)=o\), and set
\[
 \alpha(r)=R(r,0)\in L_1,
 \qquad \beta(r)=R(0,r)\in L_2,
 \qquad \gamma(r)=R(r,r)\in X.
\]
There are mutually inverse local biholomorphism germs \(H_r,K_r\)
at \((\alpha(r),\beta(r))\) and \(\gamma(r)\), respectively,
which give analytic continuation of \(H_o\) along \((\alpha,\beta)\)
and of \(K_o\) along \(\gamma\).
\end{proposition}

\begin{proof}
Transport a first-coordinate germ from the \(L_1\)-plaque at
\(\alpha(r)\) along the vertical path
\(u\mapsto R(r,ru)\), \(0\leq u\leq1\). To define this transport,
subdivide the path into finitely many subsegments lying in product
charts, and choose a local separated transition at each junction.
Within a subsegment the first coordinate is constant; at a junction it
changes by the first-coordinate local biholomorphism of that transition.
Their composition transports a variable first coordinate to the
endpoint. Refining the subdivision or changing intermediate charts
does not change the germ: separated transitions satisfy the cocycle
identity. On a common vertical subsegment, the first-coordinate transition
germ at the fixed first coordinate is locally constant, since local
separated transitions are independent of the second coordinate. It is
therefore constant along the connected subsegment. Common refinements identify any two
constructions. Similarly, transport a second-coordinate germ from the
\(L_2\)-plaque at \(\beta(r)\) along
\(u\mapsto R(ru,r)\). Combining the two transported coordinates in
one product chart at \(\gamma(r)\) gives a local biholomorphism germ
\(H_r\); let \(K_r\) be its inverse.

We must show genuine analytic continuation, rather than merely
continuous dependence of a collection of germs. Fix \(r_0\). Choose
the finite chart chains for the two reference transport paths, a fixed
small overlap neighborhood at each junction, and a common final chart
\(z\) around \(\gamma(r_0)\). By joint continuity of \(R\), the
same compact subsegments lie in the same charts, and the same junctions
lie in the same overlap neighborhoods, for all \(r\) sufficiently
near \(r_0\). Also \(\alpha(r)\) and \(\beta(r)\) remain in
fixed intrinsic leaf plaques by Lemma~\ref{lem:leaves}.

Use coordinates \(a\) and \(b\) on those two starting plaques.
Every first-coordinate transition in the vertical chain is now one
fixed holomorphic map, independent of the second coordinate. After
shrinking the initial coordinate neighborhood, their finite composition
is a single local biholomorphism \(A\) defined near
\(a(\alpha(r_0))\). It describes the transport for every nearby
\(r\), not just for \(r_0\). The horizontal chain similarly gives
one fixed second-coordinate local biholomorphism \(B\). Consequently
one holomorphic map
\begin{equation}\label{eq:fixed-transport}
 (x,y)\longmapsto
 z^{-1}\bigl(A(a(x)),B(b(y))\bigr)
\end{equation}
represents \(H_r\) at \((\alpha(r),\beta(r))\) for every nearby
\(r\), and sends that point to \(\gamma(r)\). After one further
shrinking, its single holomorphic inverse represents all the nearby
\(K_r\). At \(r=0\), both transport paths are constant, so these
germs are precisely \(H_o\) and \(K_o\). This is the required
analytic continuation.
\end{proof}

\begin{proof}[Proof of Theorem~\ref{thm:main}]
Every path from \((o,o)\) in \(L_1\times L_2\) consists of two
intrinsically continuous leaf paths. Proposition~\ref{prop:rectangles}
provides a box with these as its axes, and
Proposition~\ref{prop:germ-transport} continues \(H_o\) along the
given path. Conversely, every path from \(o\) in \(X\) is the
diagonal of a box, so the same two Propositions continue \(K_o\)
along every such path.

Let \(\pi:\widetilde X\to X\) and
\(q_i:Y_i=\widetilde L_i\to L_i\) be the based connected simply
connected covers, and put \(q=q_1\times q_2\). The product
\(Y_1\times Y_2\) is the universal cover of \(L_1\times L_2\).
The monodromy theorem, applied on these simply connected source covers,
gives holomorphic maps
\[
 H:Y_1\times Y_2\longrightarrow X,
 \qquad K:\widetilde X\longrightarrow L_1\times L_2
\]
with the prescribed initial germs. Here the usual monodromy argument
applies to manifold-valued holomorphic germs: uniqueness follows from
the identity theorem, and finite chains of local representatives are
stable under small path variations, so homotopic paths with fixed
endpoints produce the same terminal germ.

Since both source manifolds are simply connected, these maps lift to
holomorphic maps
\[
 \widetilde H:Y_1\times Y_2\longrightarrow\widetilde X,
 \qquad
 \widetilde K:\widetilde X\longrightarrow Y_1\times Y_2.
\]
Choose the lifts that send the based points to one another. Their germs
there are mutually inverse, because \(H_o\) and \(K_o\) were inverse.
Thus \(\widetilde H\circ\widetilde K\) and
\(\widetilde K\circ\widetilde H\) agree with the respective identity
maps on neighborhoods of the basepoints. The identity theorem on the
connected source manifolds makes both compositions identities
everywhere. In particular, \(\Phi=\widetilde K\) is a biholomorphism
\(\widetilde X\simeq Y_1\times Y_2\).

It remains to retain the specified splitting. Near the basepoint, the
initial product chart gives
\[
 \dd\Phi(\pi^*E_1)\subset\pr_1^*T_{Y_1},\qquad
 \dd\Phi(\pi^*E_2)\subset\pr_2^*T_{Y_2}.
\]
Equivalently, the two wrong-factor components of \(\dd\Phi\) vanish
there. They are holomorphic bundle maps, so they vanish on all of the
connected manifold \(\widetilde X\). Finally, \(\dd\Phi\) is
invertible and \(\dim_\C Y_i=r_i\) by Lemma~\ref{lem:leaves}; hence
these inclusions are equalities. The factors are connected and simply
connected by construction, proving the asserted compatible splitting.
\end{proof}

\section{Nef and big canonical bundle: parallelism and integrability}
\label{sec:canonical}

Throughout this section, \(X\) and \(T_X=E_1\oplus E_2\) satisfy the
hypotheses of Corollary~\ref{thm:nef-big}; integrability is not assumed.
We prove it by an estimate for the second fundamental forms of the
summands. This estimate also gives parallelism on the regular canonical
locus, the additional input needed for the alternative proof in the
appendices.

\begin{proposition}\label{canon:contraction}
There is a projective birational morphism with connected fibers
\[
 f:X\longrightarrow Z
\]
onto a normal projective variety with canonical singularities and ample
\(\mathbf Q\)-Cartier canonical divisor.  Moreover,
\(K_X=f^*K_Z\), and \(f\) is an isomorphism over \(Z_{\mathrm{reg}}\).
\end{proposition}

\begin{proof}
The base-point-free theorem applies to \(D=K_X\): \(D\) is nef and
\(aD-K_X=(a-1)K_X\) is nef and big for an integer \(a>1\)
\cite[Theorem~3.3]{KM98}.  The associated contraction, with its Stein
factorization, has normal projective target and connected fibers.
For some positive integer \(m\) and ample line bundle \(L\) on \(Z\),
\[
 \mathcal O_X(mK_X)\simeq f^*L.
\]
Bigness makes \(f\) generically finite; connected fibers then make it
birational.  A proper birational morphism onto a normal variety is an
isomorphism over an open set whose complement has codimension at least
two.  On this open set \(L\) agrees with
\(\mathcal O_Z(mK_Z)\).  Reflexive extension therefore identifies
\(L\) with \(\mathcal O_Z(mK_Z)\) everywhere.  Thus \(K_Z\) is
\(\mathbf Q\)-Cartier and ample.

Choose compatible canonical divisors.  The divisor
\(K_X-f^*K_Z\) is exceptional and \(\mathbf Q\)-linearly trivial.
A principal exceptional divisor is zero: the corresponding rational
function has no divisor on the normal projective variety \(Z\), hence
is constant.  Consequently \(K_X=f^*K_Z\).  The discrepancies on any
higher resolution are nonnegative because \(X\) is smooth, so \(Z\)
has canonical, and in particular klt, singularities.

Over \(Z_{\mathrm{reg}}\), the relative Jacobian has zero divisor.
It is therefore nowhere zero, and \(f\) is a local biholomorphism
there.  Properness and birationality imply that this restriction is an
isomorphism.
\end{proof}

Set
\[
 X^0=f^{-1}(Z_{\mathrm{reg}}).
\]
This is a connected dense open subset of \(X\); we identify it with
\(Z_{\mathrm{reg}}\).  Our convention is
\(\Ric(\omega)=-\sqrt{-1}\,\partial\bar\partial\log\omega^n\),
so \([\Ric(\omega)]=2\pi c_1(X)\).

\begin{lemma}[The canonical metric]\label{canon:metric-input}
Fix a Kähler form \(\omega_A\) from a polarization of \(X\).
There are smooth Kähler metrics \(g_\varepsilon\) on \(X\), for
\(0<\varepsilon<1\), whose forms satisfy
\[
 \Ric(\omega_\varepsilon)
       =-\omega_\varepsilon+\varepsilon\omega_A,
 \qquad
 [\omega_\varepsilon]
       =2\pi c_1(K_X)+\varepsilon[\omega_A].
\]
As \(\varepsilon\to0\), they converge smoothly on compact subsets of
\(X^0\) to a Kähler metric \(g\) with \(\Ric(g)=-g\).
The completion of the intrinsic Riemannian distance of
\((X^0,g)\) is naturally homeomorphic to \(Z\).
\end{lemma}

\begin{proof}
Proposition~\ref{canon:contraction} makes \(f\) a smooth projective
crepant resolution of a canonical model.  The completion statement is
\cite[Theorem~1.2]{Song14}.  The approximating metrics and their smooth
convergence are \cite[Equation~(4.28) and Lemma~4.1]{Song14}, with
\(\varepsilon=e^{-t}\).  The cohomology identity follows directly from
the Ricci equation.
\end{proof}

\begin{proposition}\label{canon:parallel}
The connection of \(g\) preserves both \(E_1|_{X^0}\) and
\(E_2|_{X^0}\).  Both subbundles are integrable on all of \(X\).
They admit simultaneous local holomorphic product coordinates.
\end{proposition}

\begin{proof}
We use the subbundle-curvature method of approximate
K\"ahler--Einstein geometry, developed in the stability results of
Enoki \cite{Enoki88} and used for parallel splittings by Guenancia
\cite[Section~3.4, Step~4, Equations~(3.12)--(3.13)]{Guenancia16}.
Here the equality of the two summands' total first Chern class with
$c_1(T_X)$ gives a direct cancellation.
Give \(T_X\) the Hermitian metric induced by \(g_\varepsilon\).
Let \(P_{j,\varepsilon}\) be orthogonal projection onto \(E_j\),
and let \(B_{j,\varepsilon}\) be its second fundamental form.
These orthogonal projections need not be projections along the
specified complementary summands.

Write \(F_{V}\) for the Chern curvature of a Hermitian bundle \(V\).
The subbundle curvature formula gives
\[
 \tr\!\left(\sqrt{-1}\Lambda_{\omega_\varepsilon}F_{E_j}\right)
 =
 \tr\!\left(
   P_{j,\varepsilon}\sqrt{-1}\Lambda_{\omega_\varepsilon}F_{T_X}
 \right)-|B_{j,\varepsilon}|^2.
\]
Here and below the second fundamental forms have their complex
Hermitian norms.  This formula follows by writing the connection in
blocks for \(E_j\oplus E_j^\perp\); the off-diagonal blocks are
\(B_{j,\varepsilon}\) and \(-B_{j,\varepsilon}^*\).
In the Kähler case the mean curvature endomorphism
\(\sqrt{-1}\Lambda_{\omega_\varepsilon}F_{T_X}\) is the Ricci
endomorphism
\[
 H_\varepsilon=-I+\varepsilon A_\varepsilon,
 \qquad A_\varepsilon=\omega_\varepsilon^{-1}\omega_A>0.
\]

Integrating the trace formula computes first Chern degrees.
Since \(c_1(E_1)+c_1(E_2)=c_1(T_X)\), and writing
\(dV_\varepsilon=\omega_\varepsilon^n/n!\), we obtain
\[
 \int_X\sum_{j=1}^2|B_{j,\varepsilon}|^2\,dV_\varepsilon
 =
 \int_X\tr(C_\varepsilon H_\varepsilon)\,dV_\varepsilon,
 \qquad
 C_\varepsilon=P_{1,\varepsilon}+P_{2,\varepsilon}-I.
\]
The ranks add to \(n\), so \(\tr C_\varepsilon=0\).
Moreover \(-I\leq C_\varepsilon\leq I\), since each
\(P_{j,\varepsilon}\) is an orthogonal projection.  Positivity of
\(A_\varepsilon\) therefore gives
\begin{align}
 0
 &\leq\int_X\sum_j|B_{j,\varepsilon}|^2\,dV_\varepsilon
 \leq\varepsilon\int_X\tr(A_\varepsilon)\,dV_\varepsilon
 \notag\\
 &=\frac{\varepsilon}{(n-1)!}
       \int_X\omega_A\wedge\omega_\varepsilon^{n-1}
   =O(\varepsilon).
 \label{canon:energy}
\end{align}
The last integral is a bounded polynomial in \(\varepsilon\), by
Lemma~\ref{canon:metric-input}.  No angle bound between the summands
is used.

On compact subsets of \(X^0\), the metrics, their connections, the
orthogonal projections, and the second fundamental forms converge
smoothly.  A nonzero limiting second fundamental form would have a
fixed positive integral on some compact neighborhood, contradicting
\eqref{canon:energy}.  Both limiting second fundamental forms are
therefore zero.

The connection of a Kähler metric is torsion-free.  Hence each
summand is integrable on \(X^0\).  Its Frobenius obstruction
\[
 \Lambda^2E_j\longrightarrow T_X/E_j,\qquad
 u\wedge v\longmapsto [u,v]\bmod E_j,
\]
is a holomorphic bundle morphism on \(X\).  It vanishes on the
dense open set \(X^0\), hence everywhere.
Local Frobenius submersions for the two complementary foliations,
taken together, give holomorphic coordinates
\((z,w)\) in which the summands are the \(z\)- and \(w\)-directions.
Transitions between such charts locally preserve the two factors.
\end{proof}


\begin{proof}[Proof of Corollary~\ref{thm:nef-big}]
A smooth projective complex manifold is compact K\"ahler. By
Proposition~\ref{canon:parallel}, the supplied summands are integrable.
Theorem~\ref{thm:main} therefore applies and gives precisely the compatible
product claimed in the corollary.
\end{proof}

The argument just given uses the canonical metric only to establish
integrability. The appendices retain a second route to the global
product: the metric completion controls the size of crossings near an
entire exceptional fibre, while rational curves provide their holomorphic
continuation. The parallelism proved here is needed for that metric
argument, even though integrability alone also follows from
\cite[Theorem~1.3]{Horing07}.

\appendix
\section{The canonical distance and its fibres}\label{app:collapsed}

This appendix begins an alternative proof of the nef-and-big application,
Corollary~\ref{thm:nef-big}. It uses the canonical metric and the geometry of
its collapsed fibres, rather than the holomorphic box-extension theorem
used in the proof of Theorem~\ref{thm:main}. The three subsequent
appendices construct crossings around each whole fibre, prove that these
crossings do not increase the canonical distance in either variable,
and use that estimate to fill finite grids along arbitrary paths.

Let \(X\) satisfy the hypotheses of Corollary~\ref{thm:nef-big}.
Proposition~\ref{canon:contraction} gives the crepant canonical morphism
\(f:X\to Z\), and Lemma~\ref{canon:metric-input} gives its limiting
K\"ahler--Einstein metric \(g\) on
\(X^0=f^{-1}(Z_{\reg})\). By Proposition~\ref{canon:parallel}, the
supplied summands are parallel on \(X^0\) and integrable on all of
\(X\). We first make them orthogonal without changing the completion,
then establish the two topological properties of the fibres that the
crossing argument needs.

\begin{proposition}\label{canon:pseudodistance}
There is a real-analytic Kähler metric \(h\) on \(X^0\), uniformly
equivalent to \(g\), for which the given splitting is orthogonal and
locally a Riemannian product.  Its intrinsic distance completion is
canonically homeomorphic to \(Z\).  Denoting the resulting distance
on \(Z\) by \(\rho\), the function
\begin{equation}\label{canon:d}
 d(x,y)=\rho(f(x),f(y))
\end{equation}
is a continuous pseudodistance on \(X\), and \(d(x,y)=0\) exactly
when \(f(x)=f(y)\).
\end{proposition}

\begin{proof}
Let \(Q_j:T_X\to E_j\) denote projection along the other supplied
summand, and define
\[
 h(v,w)=g(Q_1v,Q_1w)+g(Q_2v,Q_2w).
\]
Proposition~\ref{canon:parallel} implies
\(\nabla^gQ_j=0\), hence \(\nabla^gh=0\).
Thus \(\nabla^g\) is also the Levi-Civita connection of \(h\).
It preserves the complex structure, so \(h\) is Kähler.

The positive \(g\)-self-adjoint endomorphism representing \(h\)
is parallel.  Its eigenvalues are constant on connected \(X^0\).
Consequently there are global positive constants \(a,b\) such that
\[
 ag\leq h\leq bg.
\]
In local product coordinates, first- and second-factor real coordinate
fields \(V,W\) commute.  The equality
\(\nabla_VW=\nabla_WV\), together with preservation of the
complementary summands, makes both sides zero.  Since \(\nabla h=0\),
the first block of \(h\) depends only on first coordinates and the
second block only on second coordinates.

A Kähler--Einstein metric is real analytic in holomorphic coordinates
and their conjugates \cite[Theorem~6.1(c)]{DK81}.  The \(Q_j\) are
holomorphic, so the displayed formula for \(h\) proves its real
analyticity.

Uniform equivalence of the metrics makes the identity on \(X^0\)
bi-Lipschitz for their intrinsic distances.  It therefore extends
uniquely to a bi-Lipschitz homeomorphism of their completions.
Lemma~\ref{canon:metric-input} identifies that completion with \(Z\),
extending the given identification on \(X^0\).  The assertions about
\(d\) follow from continuity of \(f\) and from \(\rho\) being a
distance.
\end{proof}

\begin{lemma}\label{canon:fibers}
For every \(z\in Z\), the reduced fiber
\(P=f^{-1}(z)_{\mathrm{red}}\) is a compact connected projective
algebraic set.  Any two points of \(P\) are joined inside \(P\) by a
chain of images of rational curves, and \(\pi_1(P)=1\) in the ordinary
complex topology.
\end{lemma}

\begin{proof}
Projectivity and connectedness follow from the construction of \(f\).
The rational-chain assertion is the fiber theorem for resolutions of
klt varieties \cite[Corollary~1.6, arXiv v2]{HM07}.

For simple connectedness we use the resolution-fiber consequence of
Takayama's local argument: the fiber of a proper birational morphism
from a smooth complex variety to a klt variety is simply connected in
the ordinary topology
\cite[Theorem~1.1 and the first paragraph on p.~829]{Takayama03}.
This consequence is also used explicitly in
\cite[proof of Theorem~3.10, arXiv v4]{NS24}.
Here \(X\) is smooth and \(Z\) is klt by
Proposition~\ref{canon:contraction}, so it applies to \(f:X\to Z\)
and gives \(\pi_1(f^{-1}(z))=1\).
Replacing the fiber by its reduction does not change its underlying
topological space, hence \(\pi_1(P)=1\).
\end{proof}

The two conclusions of Lemma~\ref{canon:fibers} have distinct roles.
Rational chains will supply continuation routes, whereas ordinary simple
connectedness will remove their monodromy. The latter conclusion is not
inferred from rational chain connectedness of a possibly singular fibre.

\section{Rational boxes and ambient crossing germs}\label{app:rational}

Our immediate objective is to continue ordinary local crossing along
rational trees, with enough information to retain a holomorphic germ on
the ambient manifold at every parameter pair. Throughout this appendix,
$M$ is a smooth projective complex manifold with
a holomorphic splitting
\[
T_M=E_1\oplus E_2
\]
into integrable subbundles. Write $\mathcal F_i$ for the corresponding
regular foliations. Frobenius submersions for the two foliations, taken
together, give holomorphic product charts $z=(z_1,z_2)$ in which $E_i$
is the $i$th coordinate distribution. Changes between these charts have,
locally, the form
\[
(z_1,z_2)\longmapsto\bigl(F(z_1),G(z_2)\bigr).
\]
Consequently the local crossing
\begin{equation}\label{rat:local-crossing}
\kappa(x,y)=z^{-1}\bigl(z_1(x),z_2(y)\bigr)
\end{equation}
is independent, as a germ along the diagonal, of the product chart.
Our eventual application, Theorem~\ref{thm:crossing}, requires continuation
of these \emph{ambient} germs, rather than just crossing values on curves.

\begin{definition}\label{rat:definition}
Let $A,B$ be locally path connected spaces. A continuous map
$R:A\times B\to M$ is a \emph{box} if, near every parameter pair, there
are a product neighborhood $A'\times B'$ and a product chart on its
image in which
\[
z\bigl(R(a,b)\bigr)=\bigl(r_1(a),r_2(b)\bigr).
\]
Thus first-parameter slices locally follow $\mathcal F_1$ plaques and
second-parameter slices locally follow $\mathcal F_2$ plaques. In
particular these slices are continuous into their leaves with the leaf
topology. A \emph{rational tree} is a connected nodal curve whose
irreducible components are copies of $\Pp^1$ and whose dual graph is a
finite tree.
\end{definition}

For continuous maps on products of rational trees, holomorphic on each
product of components, the box condition follows from the corresponding
separate tangencies. Indeed, in a product chart each wrong-coordinate
function is constant on every branch of its parameter slice; continuity
makes these constants agree at a node.

The local charts agree with those of Section~\ref{sec:prelim}; the
definition here also allows singular and merely topological parameter
spaces. We recall two elementary facts about holonomy in this setting. A
continuous plaque-wise path in an $\mathcal F_2$ leaf transports first
transverse coordinates by composing the first-coordinate changes in a
finite chain of product charts. This gives a biholomorphic germ between
first transversals. Reversing the path gives the inverse germ. The germ
is unchanged under a leafwise homotopy with fixed endpoints: subdivide
the compact homotopy square into pieces mapping into foliation charts,
and cancel the transverse changes around their boundaries. The same
statements hold with the two foliations interchanged.

\begin{lemma}[Germs supplied by a box]\label{rat:germs}
Let $T$ be a rational tree or an interval, let $R:T\times T\to M$ be a
box, and put $u(s)=R(s,s)$. There is a coherent family of holomorphic
map germs
\[
c_{s,t}:(M\times M,(u(s),u(t)))\longrightarrow(M,R(s,t))
\]
which agrees with ordinary crossing when $s=t$ and satisfies
\begin{equation}\label{rat:differential}
\begin{aligned}
\dd_1c_{s,t}(T_M)&\subset E_1,&\qquad \dd_1c_{s,t}(E_2)&=0,\\
\dd_2c_{s,t}(T_M)&\subset E_2,&\qquad \dd_2c_{s,t}(E_1)&=0.
\end{aligned}
\end{equation}
Here coherence means that near each parameter pair a single ambient
holomorphic map represents all the nearby germs. This family is unique
among coherent families agreeing with ordinary crossing on the diagonal.
\end{lemma}

\begin{proof}
Fix $s,t$ and a path $\gamma$ in $T$ from $s$ to $t$. Set
$p=u(s)$, $q=u(t)$ and $r=R(s,t)$, and choose product charts
$z,\zeta,w$ at $p,q,r$, respectively. The path $R(s,\gamma)$ runs
from $p$ to $r$ in an $\mathcal F_2$ leaf. Its holonomy gives a
first-coordinate germ $h_1:z_1\to w_1$. Similarly,
$R(\gamma^{-1},t)$ gives a second-coordinate germ
$h_2:\zeta_2\to w_2$. Define
\begin{equation}\label{rat:germ-formula}
c_{s,t}(x,y)=w^{-1}\bigl(h_1(z_1(x)),h_2(\zeta_2(y))\bigr).
\end{equation}
Chart changes cancel in this formula, and
\eqref{rat:differential} follows immediately. Both slices are continuous
into their leaves. Since $T$ is simply connected, two choices of
$\gamma$ give fixed-endpoint leafwise homotopies, so the formula is
independent of this choice.

For completeness, coherence concerns the entire formula, not merely its
value. Choose finite chains of product charts along
$R(s,\gamma)$ and $R(\gamma^{-1},t)$ realizing $h_1,h_2$. For $s'$ near
$s$ and $t'$ near $t$, continuity on these compact paths ensures that
the same chains apply to $R(s',\gamma)$ and $R(\gamma^{-1},t')$. Join
$s'$ to $s$ and $t$ to $t'$ inside small connected parameter
neighborhoods. All the resulting endpoint segments lie in the endpoint
product charts, so introduce no additional transverse-coordinate
change. The same holomorphic representative
\eqref{rat:germ-formula} therefore represents $c_{s',t'}$. At a node
one uses a small star neighborhood and connectors through the node;
no differentiability of the parameters is involved.

When $s=t$, choose the constant path to recover
\eqref{rat:local-crossing}. Finally, a coherent family is analytic
continuation along paths in the connected space $T\times T$ of its
germ at any chosen diagonal parameter. Uniqueness of continuation of
ambient holomorphic germs proves the asserted uniqueness.
\end{proof}

\begin{remark}\label{rat:coherence}
Coherence in Lemma~\ref{rat:germs} is local over the parameter space.
Distinct parameter pairs with the same ambient input need not yet give
the same germ. That identification is the purpose of the monodromy
argument in Theorem~\ref{thm:crossing}. Likewise, the value and
\eqref{rat:differential} alone do not specify an ambient germ.
\end{remark}

\begin{lemma}[Extension of a rational box]\label{rat:extension}
Put $S=\Pp^1\times\Pp^1$. A holomorphic box on a neighborhood of the
diagonal, or of the union
\[
D=(\Pp^1\times\{t_0\})\cup(\{s_0\}\times\Pp^1),
\]
extends uniquely to a holomorphic box on $S$.
\end{lemma}

\begin{proof}
In either case the divisor has class $(1,1)$ and is a hyperplane
section of the Segre embedding of $S$. Its complement $V$ is therefore
a smooth affine, hence Stein, surface. Shrink the given neighborhood
to a connected neighborhood $W$ of the divisor. Removing a complex
hypersurface does not disconnect $W$, so $W\setminus D$ is connected,
where here $D$ also denotes the diagonal in the first case. The set
$K=S\setminus W$ is compact in $V$, and $V\setminus K=W\setminus D$.
Meromorphic Hartogs extension on Stein surfaces, in the form of
\cite[Theorem~2.2]{MP09}, extends meromorphic functions from
$V\setminus K$ to $V$.

Embed $M$ in projective space and choose a coordinate hyperplane not
containing the image of the given map. Extend its projective coordinate
ratios and glue to their original values on $W$. They define a
meromorphic map $F:S\dashrightarrow M$: the equations of the embedded
$M$ hold on $W$ and hence identically. The two slice tangencies also
persist on the holomorphic locus by the identity theorem.

The meromorphic graph is a closed analytic subset of the projective
manifold $S\times M$, so it is algebraic by Chow's theorem
\cite[Proposition~13]{Serre56}. Thus $F$ is a rational map. Clearing
the poles of its finitely many projective coordinate ratios expresses
it by a finite-dimensional linear system on $S$. Remove the fixed
divisorial part of this system and resolve its base ideal by point
blowups. This is the classical surface elimination-of-indeterminacy argument
\cite{Beauville96}; the local intersection calculation gives termination
in the present setting.
Successively blow up the remaining basepoints: the local intersection number of two general
members of the resulting linear system drops by their positive
multiplicity product at each unresolved basepoint. The finite sum of
these intersection numbers therefore forces termination. This gives
$\mu:\widehat S\to S$ and a holomorphic map $H:\widehat S\to M$.
We show that every exceptional curve is contracted by $H$. At a generic
point of an exceptional component $C$ over $(s_0,t_0)$, choose
coordinates $(x,y)$ with $C=\{x=0\}$. For local product coordinates
$s,t$ downstairs, write
\[
s\circ\mu-s_0=x^a U(x,y),\qquad
t\circ\mu-t_0=x^b V(x,y),
\]
where $a,b>0$ and $U,V$ are units. The holomorphic vector fields
\[
W_s=-xU_y\partial_x+(aU+xU_x)\partial_y,
\qquad
W_t=-xV_y\partial_x+(bV+xV_x)\partial_y
\]
annihilate $\dd(s\circ\mu)$ and $\dd(t\circ\mu)$, respectively.
Off the exceptional locus, slice tangency gives
\[
\dd H(W_s)\in E_2,\qquad \dd H(W_t)\in E_1.
\]
These inclusions extend holomorphically to $C$. There both vector
fields are nonzero multiples of $\partial_y$, whence
\[
\dd H(T_C)\subset E_1\cap E_2=0.
\]
Thus $H$ is constant on each exceptional component. Fibers of $\mu$
are connected, so each whole fiber has a single image. Around this
image choose a target coordinate chart; properness gives a neighborhood
downstairs whose inverse image maps into the chart. Its coordinate
functions descend holomorphically, since
$\mu_*\mathcal O_{\widehat S}=\mathcal O_S$, equivalently by removal
of the isolated singularity downstairs. Hence $F$ is holomorphic.
Its slice tangencies give the box condition, and uniqueness follows
from the identity theorem.
\end{proof}

We have extended a box from one diagonal or one pair of axes. A rational
tree has a unique route between components; this lets us repeat the axes
construction without making a new compatibility choice.

\begin{proposition}[Rational boxes]\label{rat:boxes}
For every holomorphic map $u:T\to M$ from a rational tree there is a
unique continuous box
\[
R_u:T\times T\longrightarrow M,
\qquad R_u(s,s)=u(s),
\]
holomorphic on every product of irreducible components. If a rational
tree $T'$ contains $T$ as a subtree and $u':T'\to M$ extends $u$, then
$R_{u'}|_{T\times T}=R_u$. The ambient germ families of
Lemma~\ref{rat:germs} agree on this restriction as well.
\end{proposition}

\begin{proof}
First take $T=\Pp^1$. Applying local crossing to $u(s),u(t)$ gives
coherent holomorphic germs along the diagonal of $S$. Compactness
allows them to be glued on a neighborhood of the diagonal: take a
finite family of representatives, shrink their domains, and then
shrink around the diagonal so that overlapping representatives agree.
Lemma~\ref{rat:extension} gives a box on $S$. In a product chart any
box near $(s,s)$ has the form $(a(s),b(t))$; its diagonal specifies
both functions. This proves uniqueness near the diagonal and hence
on $S$.

We also need the corresponding construction from axes. Suppose
$\alpha:\Pp^1\to M$ is tangent to $E_1$,
$\beta:\Pp^1\to M$ is tangent to $E_2$, and
$\alpha(s_0)=\beta(t_0)$. Transport the germ of $\beta$ at $t_0$
along $\alpha$, using $\mathcal F_1$ holonomy. There is no path
ambiguity: a loop in the parameter sphere contracts, and its image
gives a leafwise homotopy. In product charts the resulting map near
$(s,t_0)$ has first coordinate supplied by $\alpha(s)$ and second
coordinate a holomorphic germ in $t$. It is jointly holomorphic and
the germs are coherent by the finite-chart argument above. Do the
symmetric construction along the other axis. Near their intersection
both maps are ordinary crossing, so, after shrinking, they glue on a
neighborhood of the axes. Lemma~\ref{rat:extension} extends the box
over $S$. The two prescribed axes determine its germ at $(s_0,t_0)$,
so the extension is unique. This includes constant curves.

For a general tree, first construct the box on the square of each
component. Proceed through ordered pairs of distinct components by
increasing distance in the dual tree. On $C_i\times C_j$, prescribe
the two axes at the nodes pointing toward each other along the unique
component route. These axes have already been defined at smaller
distances; they are tangent to $E_1,E_2$, respectively, and agree at
their corner by the previous edge identifications. Apply the axes
construction. Every common edge matches: for two adjacent components,
their distances to a fixed third component differ by one, and the
common edge is an inward axis for the later pair. The finitely many
component maps therefore glue continuously, and their separate
tangencies imply the box condition also at nodes.

The same induction proves uniqueness, first on component squares and
then from the inward axes. Restriction to a subtree is consequently
$R_u$. To compare its ambient germs, compute the two holonomies along
paths lying in $T$. The paths and box values are unchanged in $T'$,
so Equation~\eqref{rat:germ-formula} gives precisely the same germs.
\end{proof}

\section{Crossing near a rationally chain connected set}\label{app:crossing}

The rational boxes now give a crossing map on a neighborhood of an
entire projective set. The passage has two tasks: provide continuation
along enough paths in the set, and remove path dependence while retaining
ambient holomorphic germs. In this appendix \(M\) is smooth and projective, and
\(T_M=E_1\oplus E_2\) is an integrable splitting. A set is
\emph{joined by rational chains} if every two of its points can be
joined inside it by a finite chain of images of \(\Pp^1\).

\begin{theorem}[Neighborhood crossing]\label{thm:crossing}
Let \(P\subset M\) be a nonempty connected projective algebraic set.
Suppose that \(P\) is simply connected and joined by rational chains.
There are a connected open neighborhood \(U\) of \(P\) and a holomorphic
map \(c:U\times U\to M\) that agrees near the diagonal with ordinary
crossing in product charts. Throughout \(U\times U\), it satisfies
\begin{equation}\label{cross:identities}
\begin{gathered}
 c(x,x)=x,\qquad
 \dd_1c(T_M)\subset E_1,\quad \dd_1c(E_2)=0,\\
 \dd_2c(T_M)\subset E_2,\quad \dd_2c(E_1)=0.
\end{gathered}
\end{equation}
The image subbundles in these formulas are evaluated at \(c(x,y)\).
\end{theorem}

\subsection{Continuation of ambient germs}

Put \(Q=P\times P\), fix \(p\in P\), and write \(q_*=(p,p)\).
All germs below are germs of holomorphic maps from \(M\times M\) to \(M\),
even when their basepoints lie in \(Q\).

\begin{definition}\label{cross:attained}
A germ over \((a,b)\in Q\) is \emph{attained} if it is the crossing germ
of Lemma~\ref{rat:germs} associated to a rational-tree box \(R_u\) of
Proposition~\ref{rat:boxes}, at two distinct smooth marked points \(s,t\) of
its parameter tree, with \(u(s)=a\) and \(u(t)=b\).
\end{definition}

The ordinary crossing at \(q_*\) is attained using a constant sphere.
We do not assume that different realizing trees give the same germ.

\begin{lemma}\label{cross:tree-path}
Let \(k:D\to Q\) be a holomorphic map from a rational tree with distinct
smooth markings \(m_1,m_2\). Every attained germ at \(k(m_1)\) continues
along the image of any path from \(m_1\) to \(m_2\). Its terminal germ
is attained and is independent of the path in \(D\).
\end{lemma}
\begin{proof}
Choose a realizing tree \(T\) with markings \(s,t\) for the initial
germ. Attach two separate copies of \(D\) to \(T\), identifying their
first markings with \(s\) and \(t\), respectively. Map the copies to
\(M\) by the two components of \(k\). The resulting curve is again a
rational tree, and the maps agree at the attachment points. Its box
restricts to the original box, including its ambient germs, by the
uniqueness and branch compatibility in Proposition~\ref{rat:boxes}
and Lemma~\ref{rat:germs}.

Move the two parameter points along the copies of a path in \(D\).
Lemma~\ref{rat:germs} gives exactly the asserted continuation along
its image under \(k\). Its terminal markings are smooth and distinct,
so the terminal germ is attained. All paths between the markings are
homotopic in the simply connected tree; the corresponding leafwise
holonomy is unchanged.
\end{proof}

We will use the following elementary stability fact. It applies to
holomorphic germs on any complex manifold, with values in another complex
manifold.

\begin{lemma}[Stability of continuation]\label{cross:stability}
Suppose that a germ continues along a continuous path \(\gamma:[0,1]\to N\)
in a complex manifold \(N\). There are a representative of its terminal
germ and a uniform neighborhood of \(\gamma\) in the space of paths
starting at \(\gamma(0)\) such that continuation along every path in that
neighborhood exists and ends in that representative.
\end{lemma}
\begin{proof}
Choose a finite subdivision and holomorphic representatives on open
neighborhoods of the compact subpaths. Consecutive representatives agree
on an open neighborhood of each subdivision point. After shrinking the
subpath neighborhoods and these overlaps, every sufficiently uniformly
close path stays in the same domains and crosses through the same
equality overlaps. The same finite chain of representatives works.
\end{proof}

Continuation along a fixed path from a fixed germ is unique. Indeed,
two continuation chains that agree initially agree locally at the first
possible separation point, by the identity theorem in overlapping
ambient neighborhoods. This also proves that continuation along a path
followed by its reverse returns the initial germ.

\subsection{A compact family of rational routes}

Embed \(Q\) in a projective space \(\Pp^N\). We use the following standard
properties of genus-zero stable maps: fixed-degree spaces with three
markings have projective coarse moduli spaces; evaluations are algebraic;
maps factoring through a fixed projective subscheme form a closed locus;
and locally the projective-space moduli problem has finite-quotient charts
with universal families of nodal curves. These statements follow from
Fulton--Pandharipande \cite[Theorem~1, Proposition~3, and Sections~1.3, 4.1, 4.4, and 5.1]{FP97}.
For the local families we always use the ambient target \(\Pp^N\), which
is convex, and restrict to the locus of maps into \(Q\).

\begin{lemma}\label{cross:parameter-space}
There is a compact semialgebraic space \(S_*\) of three-pointed
rational-tree maps into \(Q\) whose first evaluation is \(q_*\) and
whose second evaluation
\[
 e:S_*\longrightarrow Q
\]
is surjective.
\end{lemma}
\begin{proof}
For each degree \(d\ge0\), take the closed locus \(S_d\) in
\(\overline M_{0,3}(\Pp^N,d)\) of maps into \(Q\) with first evaluation
\(q_*\). Its second evaluation image \(F_d\) is closed algebraic.
Every point of \(Q\) lies in some \(F_d\): join it to \(q_*\) by moving
first in one factor and then in the other along rational chains.
Represent the chain by a rational tree. Insert contracted components if
needed to put the endpoints at distinct smooth markings, add a third
marking, and stabilize. This does not change the endpoint evaluations.

Write \(Q=Q_1\cup\cdots\cup Q_r\) for its irreducible components.
For each \(i\),
\[
 Q_i=\bigcup_{d\ge0}(Q_i\cap F_d).
\]
If none of the \(F_d\) contained \(Q_i\), this would express the compact
complex variety \(Q_i\) as a countable union of closed nowhere dense
sets, contradicting the Baire category theorem. Choose \(d_i\) with
\(Q_i\subset F_{d_i}\). The disjoint union of the finitely many spaces
\(S_{d_i}\), omitting repetitions, is the required \(S_*\).
\end{proof}

\begin{lemma}[Routes in a degenerating tree]\label{cross:nearby-routes}
Fix a stable-map parameter in \(S_*\) and a path between its first two
markings, chosen along smooth component arcs and through the intervening
nodes. For every sufficiently nearby parameter one can choose a path
between its corresponding markings whose image in \(Q\), after
parametrization on \([0,1]\), is uniformly as close as desired to the
fixed image path.
\end{lemma}
\begin{proof}
Choose a representative of the limiting stable map in a local
finite-quotient chart with a family. A finite quotient map is open, so
nearby coarse parameters admit representatives close to this one;
no continuous choice is required.

Follow the finite component route between the markings. Away from
small segments around its nodes, finitely many relative smooth charts
move the chosen arcs to nearby fibers, with the prescribed marked ends.
At a node the family is locally \(xy=a\), where \(a\) is a function
of the parameter. The incoming and outgoing nearby pieces can be joined
within an arbitrarily small plumbing neighborhood when \(a\) is small;
the same assertion holds for \(a=0\). The family map is continuous,
so all images in that neighborhood lie close to the limiting nodal
value. Parametrize these joining pieces on the corresponding small
subintervals. The resulting paths have uniformly close images.
Contracted components cause no difficulty, since their limiting images
are constant. Restricting the parameter locus to maps into \(Q\)
keeps all these image paths inside \(Q\).
\end{proof}

\begin{lemma}\label{cross:semialgebraic-lift}
Every continuous semialgebraic path \(\sigma:[0,1]\to Q\) has a
finite subdivision such that on each closed subinterval it admits a
continuous lift to \(S_*\). Lifts on consecutive subintervals need not
agree at their common endpoint.
\end{lemma}
\begin{proof}
Semialgebraic choice applied to
\[
 \{(t,s)\in[0,1]\times S_*:e(s)=\sigma(t)\}\longrightarrow[0,1]
\]
gives a semialgebraic section. It is continuous on finitely many open
subintervals. Compactness of \(S_*\) and one-variable semialgebraic
monotonicity give one-sided limits at every endpoint. Continuity of
\(e\) makes those limits lift the correct endpoint values. Extend each
interval lift by its one-sided limits. We use the standard semialgebraic
choice and one-variable finiteness results in
\cite[Theorems~2.1, 3.1, and 3.4]{Coste99}; see also \cite{BCR98}.
\end{proof}

\subsection{Monodromy and an actual neighborhood}

\begin{lemma}\label{cross:all-paths}
Every attained germ continues along every continuous semialgebraic path
in \(Q\) starting at its basepoint, and the terminal germ is attained.
\end{lemma}
\begin{proof}
Use Lemma~\ref{cross:semialgebraic-lift} and consider one lifted
interval \([a,b]\), with prescribed attained germ \(g_a\) at
\(\sigma(a)\). Continue \(g_a\) backwards along a marked route in the
tree at \(a\), obtaining an attained germ \(h\) at \(q_*\) by
Lemma~\ref{cross:tree-path}. For each \(t\in[a,b]\), continue this
\emph{fixed} germ \(h\) forwards through the tree over \(t\), and denote
the endpoint germ by \(g_t\). The choice of route does not affect it.
At \(t=a\), reverse-path uniqueness gives the prescribed \(g_a\).

For \(t\) near any \(t_0\), Lemma~\ref{cross:nearby-routes} permits
uniformly close image routes from the common point \(q_*\).
Lemma~\ref{cross:stability} then supplies one ambient representative
giving all \(g_t\) near \(t_0\). Thus these germs are analytic
continuation along \(\sigma\) on the whole interval. Each is attained
by Lemma~\ref{cross:tree-path}. Repeat from the terminal germ on the
next interval. The intermediate germ \(h\) at \(q_*\) is not assumed
to be the original ordinary crossing; only its being fixed on the
current interval is used.
\end{proof}

\begin{proof}[Proof of Theorem~\ref{thm:crossing}]
The compact family of routes has supplied continuation along every
semialgebraic path, including through singularities and changes of
component. It remains to identify the terminal germs and glue them.
Triangulate \(Q\) semialgebraically
\cite[Theorem~4.4]{Coste99}. Since \(P\) is simply connected,
so is \(Q\). Start with the ordinary crossing germ at \(q_*\), and
continue it along piecewise linear paths in the triangulation. These
paths are semialgebraic in \(Q\), so Lemma~\ref{cross:all-paths}
applies. Any two such paths with the same endpoints are homotopic
relative endpoints through paths with semialgebraic slices: apply
relative piecewise linear approximation to a nullhomotopy in the
finite triangulation. For a fixed slice, compact-path stability
(Lemma~\ref{cross:stability}) and uniform continuity of the homotopy
show that nearby slices have the same terminal germ. The terminal
germ is consequently constant throughout the homotopy. We obtain
well-defined ambient germs at every point of \(Q\).

These germs are locally coherent along \(Q\). Indeed, choose a small
relative path-connected neighborhood of an endpoint inside the domain
of its representative. Appending small paths that are piecewise linear in the fixed
semialgebraic triangulation, after subdivision, and lie in that
neighborhood shows that the same representative gives the assigned
germs there. Such neighborhoods are available from the triangulation.

For completeness, this coherent family has a representative on an
actual neighborhood of \(Q\). Choose finitely many representatives
\(c_i:V_i\to M\) giving the family on \(Q\cap V_i\), and smaller open
sets \(W_i\) with compact closures in \(V_i\), still covering \(Q\).
On \(V_i\cap V_j\), the set where the germs of \(c_i,c_j\) agree is
a union of connected components, by the identity theorem. Its
complement is therefore relatively closed. The compact sets
\[
 B_{ij}=\{x\in\overline W_i\cap\overline W_j:
                 (c_i)_x\ne(c_j)_x\}
\]
are disjoint from \(Q\). Removing their finite union and restricting
to \(\bigcup_iW_i\) gives a neighborhood of \(Q\) on which the
representatives glue holomorphically. Notice that equality here is
equality of \emph{ambient germs}, not merely of values on \(Q\).

Compactness of \(P\) gives an open neighborhood \(U\) of \(P\) whose
square lies in this domain. Since \(P\) is connected, take the connected
component of \(U\) containing it. The glued map restricts to
\(c:U\times U\to M\). Its differential identities in
\eqref{cross:identities} hold near \((p,p)\) and hence throughout
connected \(U\times U\) by holomorphicity. The identity on the diagonal
extends throughout connected \(U\). In a product chart about any
\(x\in U\), these differential identities make the first output
coordinate depend only on the first coordinate of the first input,
and the second output only on the second coordinate of the second
input. The diagonal identity identifies both functions with the
corresponding coordinates. Hence \(c\) is ordinary crossing near every
diagonal point, as required.
\end{proof}

\section{Nonexpansion and finite grids}\label{app:global}

We return to the canonical morphism $f\colon X\to Z$.  Put
$X^0=f^{-1}(Z_{\reg})$, and let $h$ be the orthogonal product metric
associated with the parallel splitting on $X^0$ from
Proposition~\ref{canon:pseudodistance}.
We use the following established properties of $h$: it is real analytic,
it is locally a product metric for $E_1\oplus E_2$, and its intrinsic
length metric has a completion identified homeomorphically with $Z$.
Write $\rho$ for the resulting distance on $Z$, and recall
\[
 d(x,x')=\rho(f(x),f(x'))\qquad(x,x'\in X).
\]
Thus $d$ is a continuous pseudodistance, whose zero-distance classes are
the fibers of $f$.  Diameters below are taken for $d$.

For each reduced fiber $P=f^{-1}(z)_{\mathrm{red}}$, Theorem~\ref{thm:crossing}
provides a connected open neighborhood $U\supset P$ and a holomorphic
crossing map $c\colon U\times U\to X$.  We will use its identities
\begin{equation}\label{global:crossing-identities}
 \begin{gathered}
 c(x,x)=x,\qquad
 \dd_1c(T_X)\subset E_1,\quad \dd_1c(E_2)=0,\\
 \dd_2c(T_X)\subset E_2,\quad \dd_2c(E_1)=0.
 \end{gathered}
\end{equation}
Near a diagonal point, $c$ is the ordinary operation that takes the first
coordinate of the first input and the second coordinate of the second.

\subsection{Nonexpansion across the exceptional locus}

\begin{lemma}\label{global:nonexpansion}
For each fiber value $z\in Z$, one can choose a crossing map as above and
$r>0$ such that $f^{-1}(B_{\rho}(z,r))\subset U$ and
\begin{equation}\label{global:lipschitz}
 \begin{split}
 d\bigl(c(a,b),c(a',b)\bigr)&\leq d(a,a'),\\
 d\bigl(c(a,b),c(a,b')\bigr)&\leq d(b,b')
 \end{split}
\end{equation}
whenever all the inputs in the relevant inequality belong to
$f^{-1}(B_{\rho}(z,r/10))$.
\end{lemma}

\begin{proof}
On the open set
\[
 V=\{(a,b)\in U\times U:a,b,c(a,b)\in X^0\}
\]
we claim the tensor identities
\begin{equation}\label{global:metric-identities}
 \begin{split}
 \bigl\lVert\dd_1c(v)\bigr\rVert_{h,c(a,b)}^2
     &=\lVert v_1\rVert_{h,a}^2,\\
 \bigl\lVert\dd_2c(w)\bigr\rVert_{h,c(a,b)}^2
     &=\lVert w_2\rVert_{h,b}^2,
 \end{split}
\end{equation}
where $v=v_1+v_2$ and $w=w_1+w_2$ denote the supplied splittings of real
tangent vectors.  The identities hold near any regular diagonal point
by the local product expression for $h$.  The complement of $V$ is a
proper complex analytic subset of the connected manifold $U\times U$,
so $V$ is connected.  Both sides of each identity are real analytic on
$V$.  The real analytic identity principle proves
\eqref{global:metric-identities} throughout $V$.

Properness of $f$ and the topology induced by $\rho$ give $r>0$ with
$f^{-1}(B_{\rho}(z,r))\subset U$.  Consider the first inequality in
\eqref{global:lipschitz}.  Approximate $(a,a',b)$ in $X^3$ by triples
$(a_m,a'_m,b_m)$ for which all three inputs and both output points are in
$X^0$.  Such triples are dense: each excluded condition is a proper
analytic subset of $U^3$, and properness of the output conditions follows
from the regular values of $c$ near a regular diagonal point.

For large $m$ all inputs lie over $B_{\rho}(z,r/9)$.  Join $a_m$ to $a'_m$
in $X^0$ by a piecewise smooth path of $h$-length less than
$d(a_m,a'_m)+\eta$, with $\eta>0$ arbitrarily small.  This is possible
because $d$ on $X^0$ is its intrinsic length distance.  Every point of
this path lies over $B_{\rho}(z,r/3+\eta)$: the initial point has distance
less than $r/9$ from $z$, and the path length is less than $2r/9+\eta$.
In particular the path remains in $U$, with a margin for perturbation.

For this fixed $b_m$, the set
\[
 A_m=\{x\in U\cap X^0:c(x,b_m)\notin X^0\}
\]
is a proper analytic subset, since the output at $a_m$ is regular.
It has real codimension at least two.  Keeping the endpoints fixed, we
may perturb the path to avoid $A_m$ with arbitrarily small length error.
Indeed, on a neighborhood of the compact path the metric is smooth;
general position relative to an analytic stratification gives such a
perturbation within that neighborhood.  Along the perturbed path,
\eqref{global:metric-identities} bounds the output speed by the input
speed, because the two summands are $h$-orthogonal.  Integration, followed
by letting the length errors tend to zero, proves
\[
 d\bigl(c(a_m,b_m),c(a'_m,b_m)\bigr)\leq d(a_m,a'_m).
\]
Continuity of $c$ and $d$ gives the first inequality for the original
inputs.  Interchanging the two variables proves the second.
\end{proof}

Compactness of $Z$ now supplies a uniform constant $\delta>0$ with the
following property:
\begin{equation}\label{global:uniform}
 \begin{gathered}
 \text{Every set }S\subset X\text{ with }\operatorname{diam}_d(S)\leq2\delta
 \text{ lies in a crossing domain }U,\\
 \text{and \eqref{global:lipschitz} holds for all inputs taken from }S.
 \end{gathered}
\end{equation}
To see this, take a finite subcover of $Z$ by the smaller balls in
Lemma~\ref{global:nonexpansion}, and choose $2\delta$ strictly below a
Lebesgue number.  No ordinary diameter bound for the fibers is asserted
or needed here.

\subsection{Two finite grid constructions}

For intervals $I,J$, the box condition of
Definition~\ref{rat:definition} says that the first product coordinate of
$R:I\times J\to X$ depends locally only on the parameter in $I$, and
the second only on the parameter in $J$. Its horizontal and vertical
paths consequently follow the $E_1$- and $E_2$-plaques, respectively.
We first record the rule for filling one tile.

Suppose two prescribed edges $\alpha\colon I\to X$ and
$\beta\colon J\to X$ meet at a corner $p$, follow the first and second
foliations respectively, and each have diameter at most $\delta$.
Their union has diameter at most $2\delta$.  Choose a crossing map by
\eqref{global:uniform} and set
\begin{equation}\label{global:tile}
 R(s,t)=c(\alpha(s),\beta(t)).
\end{equation}
This is a box and agrees with both prescribed edges.  For example, for
fixed $s$, vary the second input along the part of $\alpha$ joining $p$
to $\alpha(s)$.  The identity $\dd_2c(E_1)=0$ makes the value constant,
and therefore
\[
 c(\alpha(s),p)=c(\alpha(s),\alpha(s))=\alpha(s).
\]
The other edge follows from $\dd_1c(E_2)=0$.  This argument applies to
continuous plaque paths, since the corresponding functions are locally
constant on their parameter intervals.  Finally,
\eqref{global:lipschitz} shows that every horizontal or vertical partial
slice of this tile has diameter at most $\delta$.  Thus the same bound
is available for its two new edges.

\begin{lemma}\label{global:boxes}
The following boxes exist.
\begin{enumerate}
 \item For every continuous path $\lambda\colon[0,1]\to X$, there is a
 box $R\colon[0,1]^2\to X$ such that $R(s,s)=\lambda(s)$.
 \item For continuous paths $\alpha$ and $\beta$ in leaves of the first
 and second foliations, respectively, with
 $\alpha(0)=\beta(0)$, there is a box on $[0,1]^2$ satisfying
 $R(s,0)=\alpha(s)$ and $R(0,t)=\beta(t)$.
\end{enumerate}
Here continuity of the prescribed leaf paths is in the intrinsic leaf
topology.
\end{lemma}

\begin{proof}
For the first assertion choose a subdivision
$0=t_0<\cdots<t_N=1$ such that each $\lambda(I_i)$ has diameter at most
$\delta$, where $I_i=[t_{i-1},t_i]$.  On every diagonal tile $I_i^2$ put
\[
 R(s,t)=c_i(\lambda(s),\lambda(t)),
\]
using \eqref{global:uniform}.  Its diagonal is $\lambda$, and each of
its partial slices has diameter at most $\delta$ by
\eqref{global:lipschitz}.  These tiles agree at their common vertices.

Fill the remaining tiles $I_i\times I_j$ by increasing $|i-j|$.
Above the diagonal, the bottom and right edges of a new tile have
already been defined; below it, the top and left edges have already
been defined.  In either case these are one horizontal and one vertical
edge, meeting at a corner.  Apply \eqref{global:tile}.
The induction invariant is that every partial slice of each completed
tile has diameter at most $\delta$.  The tile rule proves both agreement
on its old edges and preservation of this invariant.  Tiles at the same
stage share no edge.  If they share a vertex, its value is already fixed
by an intervening tile from an earlier stage.  Thus the finite filling
is consistent everywhere.

For the second assertion, subdivide each prescribed axis so that the
image of every edge segment has diameter at most $\delta$.  Fill the
rectangular array in increasing order of the sum of its two indices.
The bottom and left edges of a new tile have already been prescribed,
so precisely the same rule and invariant apply.

In both constructions the tile maps glue continuously.  The box
property also holds across seams: near a seam value, continuity places
the adjacent path pieces in one foliation chart, and both follow the
plaque through their common endpoint.  This proves the assertions.
\end{proof}

The filling orders are the same as in Figure~\ref{fig:grid}. Here the
one-tile rule \eqref{global:tile} propagates the diameter bound
\(\delta\) by separate nonexpansion, rather than preserving a
holomorphic parametrization on a full polydisk.

The construction uses only continuity and diameter control.  In
particular it applies to paths inside a collapsed fiber, whose
pseudodiameter is zero, without assigning those paths any finite length
for the limiting metric.

\subsection{The resulting product of universal covers}

The two grid constructions have now supplied the exact inputs for
transport of inverse germs in Proposition~\ref{prop:germ-transport}.
That proposition uses only product charts and a continuous box; it does
not use the compact K\"ahler splitting theorem or the construction of
boxes in Sections~\ref{sec:hartogs}--\ref{sec:global}. We spell out the
remaining application to keep the alternative argument independent of
that theorem.

Fix \(o\in X\), and give the leaves \(L_1,L_2\) through \(o\)
their intrinsic complex-manifold structures, as in
Lemma~\ref{lem:leaves}. Let
\[
 H_o:(L_1\times L_2,(o,o))\longrightarrow(X,o),\qquad
 K_o:(X,o)\longrightarrow(L_1\times L_2,(o,o))
\]
be inverse local product germs. Every path in \(X\) from \(o\) is
the diagonal of a box by Lemma~\ref{global:boxes}(1). Its axes lie
continuously in the fixed intrinsic leaves, so
Proposition~\ref{prop:germ-transport} continues \(K_o\) along that
path. Conversely, a path in \(L_1\times L_2\) consists of two
intrinsically continuous leaf paths. Lemma~\ref{global:boxes}(2) supplies
a box with those axes, and the same transport proposition continues
\(H_o\) along the prescribed path. All continued germs are local
biholomorphisms.

\begin{proposition}\label{global:product}
Under the hypotheses of Corollary~\ref{thm:nef-big}, let \(Y\to X\)
be the universal cover and \(\widetilde L_i\to L_i\) the universal
covers of the leaves through \(o\). Then
\[
 Y\simeq\widetilde L_1\times\widetilde L_2
\]
biholomorphically, with the lifts of \(E_i\) identified with the
respective tangent bundles of the factors.
\end{proposition}

\begin{proof}
All-path continuation and monodromy give holomorphic maps
\[
 F:Y\longrightarrow L_1\times L_2,
 \qquad G:\widetilde L_1\times\widetilde L_2\longrightarrow X.
\]
With compatible basepoints, the covering lifting theorem gives
holomorphic lifts
\[
 \widehat F:Y\longrightarrow\widetilde L_1\times\widetilde L_2,
 \qquad \widehat G:\widetilde L_1\times\widetilde L_2\longrightarrow Y.
\]
Their initial germs are inverse. Each composite is therefore the
identity near its basepoint, and hence everywhere on its connected
source by the identity theorem. The wrong-factor components of their
differentials are holomorphic bundle morphisms that vanish in the
initial product chart. They vanish identically, giving the asserted
compatibility with the supplied splitting.
\end{proof}

Proposition~\ref{global:product} completes the alternative proof of
Corollary~\ref{thm:nef-big}. Its continuation mechanism is the separate
nonexpansion of crossings around collapsed canonical fibres together
with the finite-grid diameter bounds.

\begin{thebibliography}{99}
\raggedright

\bibitem{Beauville96}
Arnaud Beauville,
\emph{Complex Algebraic Surfaces}, second edition,
London Mathematical Society Student Texts, vol.~34,
Cambridge University Press, Cambridge, 1996.
\href{https://doi.org/10.1017/CBO9780511623936}{doi:10.1017/CBO9780511623936}.

\bibitem{Beauville00}
Arnaud Beauville,
\emph{Complex manifolds with split tangent bundle},
in Thomas Peternell and Frank-Olaf Schreyer (eds.),
\emph{Complex Analysis and Algebraic Geometry: A Volume in Memory of Michael Schneider},
de Gruyter, Berlin, 2000, pp.~61--70.
\href{https://doi.org/10.1515/9783110806090-004}{doi:10.1515/9783110806090-004}.

\bibitem{BCR98}
Jacek Bochnak, Michel Coste, and Marie-Fran\c{c}oise Roy,
\emph{Real Algebraic Geometry},
Ergebnisse der Mathematik und ihrer Grenzgebiete, 3.~Folge, vol.~36,
Springer-Verlag, Berlin, 1998.
\href{https://doi.org/10.1007/978-3-662-03718-8}{doi:10.1007/978-3-662-03718-8}.

\bibitem{Brody78}
R.~Brody,
\emph{Compact manifolds and hyperbolicity},
Trans. Amer. Math. Soc. \textbf{235} (1978), 213--219.
\href{https://doi.org/10.1090/S0002-9947-1978-0470252-3}{doi:10.1090/S0002-9947-1978-0470252-3}.

\bibitem{BPT06}
Marco Brunella, Jorge Vit\'orio Pereira, and Fr\'ed\'eric Touzet,
\emph{K\"ahler manifolds with split tangent bundle},
Bulletin de la Soci\'et\'e Math\'ematique de France
\textbf{134} (2006), no.~2, 241--252.
\href{https://doi.org/10.24033/bsmf.2507}{doi:10.24033/bsmf.2507}.

\bibitem{CP11}
Fr\'ed\'eric Campana and Thomas Peternell,
\emph{Geometric stability of the cotangent bundle and the universal cover
of a projective manifold}, with an appendix by Matei Toma,
Bull. Soc. Math. France \textbf{139} (2011), no.~1, 41--74.
\href{https://doi.org/10.24033/bsmf.2599}{doi:10.24033/bsmf.2599}.

\bibitem{Coste99}
Michel Coste,
\emph{An Introduction to O-minimal Geometry},
lecture notes, November 1999.
\url{https://perso.univ-rennes1.fr/michel.coste/polyens/OMIN.pdf}.

\bibitem{Demailly02}
Jean-Pierre Demailly,
\emph{On the Frobenius integrability of certain holomorphic $p$-forms},
in \emph{Complex Geometry}, Springer, Berlin, 2002, pp.~93--98.
\href{https://doi.org/10.1007/978-3-642-56202-0_6}{doi:10.1007/978-3-642-56202-0\_6}.

\bibitem{DemaillyCADG}
J.-P.~Demailly,
\emph{Complex Analytic and Differential Geometry},
manuscript, version of June 21, 2012.
\url{https://people.math.harvard.edu/~demarco/Math274/Demailly_ComplexAnalyticDiffGeom.pdf}.

\bibitem{deRham52}
G.~de Rham,
\emph{Sur la r\'eductibilit\'e d'un espace de Riemann},
Comment. Math. Helv. \textbf{26} (1952), 328--344.
\href{https://doi.org/10.1007/BF02564308}{doi:10.1007/BF02564308}.

\bibitem{DK81}
Dennis M.~DeTurck and Jerry L.~Kazdan,
\emph{Some regularity theorems in Riemannian geometry},
Annales Scientifiques de l'\'Ecole Normale Sup\'erieure (4)
\textbf{14} (1981), no.~3, 249--260.
\href{https://doi.org/10.24033/asens.1405}{doi:10.24033/asens.1405}.

\bibitem{Druel00}
S.~Druel,
\emph{Vari\'et\'es alg\'ebriques dont le fibr\'e tangent est totalement d\'ecompos\'e},
J. Reine Angew. Math. \textbf{522} (2000), 161--171.
\href{https://arxiv.org/abs/math/9901138v2}{arXiv:math/9901138v2}.

\bibitem{DPPT22}
S.~Druel, J.~V.~Pereira, B.~Pym, and F.~Touzet,
\emph{A global Weinstein splitting theorem for holomorphic Poisson manifolds},
Geom. Topol. \textbf{26} (2022), no.~6, 2831--2853.
\href{https://doi.org/10.2140/gt.2022.26.2831}{doi:10.2140/gt.2022.26.2831}.

\bibitem{DPPT24}
St\'ephane Druel, Jorge Vit\'orio Pereira, Brent Pym, and Fr\'ed\'eric Touzet,
\emph{Numerically flat foliations and holomorphic Poisson geometry},
preprint, 2024.
\href{https://arxiv.org/abs/2411.08806v1}{arXiv:2411.08806v1}.

\bibitem{Enoki88}
Ichiro Enoki,
\emph{Stability and negativity for tangent sheaves of minimal K\"ahler spaces},
in \emph{Geometry and Analysis on Manifolds}, Lecture Notes in Mathematics,
vol.~1339, Springer, Berlin, 1988, pp.~118--126.
\href{https://doi.org/10.1007/BFb0083051}{doi:10.1007/BFb0083051}.

\bibitem{FP97}
William Fulton and Rahul Pandharipande,
\emph{Notes on stable maps and quantum cohomology},
in J\'anos Koll\'ar, Robert Lazarsfeld, and David R.~Morrison (eds.),
\emph{Algebraic Geometry---Santa Cruz 1995},
Proceedings of Symposia in Pure Mathematics, vol.~62, part~2,
American Mathematical Society, Providence, RI, 1997, pp.~45--96.
\href{https://arxiv.org/abs/alg-geom/9608011}{arXiv:alg-geom/9608011}.

\bibitem{Guenancia16}
Henri Guenancia,
\emph{Semistability of the tangent sheaf of singular varieties},
Algebraic Geometry \textbf{3} (2016), no.~5, 508--542.
\href{https://doi.org/10.14231/AG-2016-024}{doi:10.14231/AG-2016-024}.

\bibitem{HM07}
Christopher D.~Hacon and James McKernan,
\emph{On Shokurov's rational connectedness conjecture},
Duke Mathematical Journal \textbf{138} (2007), no.~1, 119--136.
\href{https://doi.org/10.1215/S0012-7094-07-13813-4}{doi:10.1215/S0012-7094-07-13813-4};
\href{https://arxiv.org/abs/math/0504330v2}{arXiv:math/0504330v2}.

\bibitem{Horing07}
Andreas H\"oring,
\emph{Uniruled varieties with split tangent bundle},
Mathematische Zeitschrift \textbf{256} (2007), no.~3, 465--479.
\href{https://doi.org/10.1007/s00209-006-0072-5}{doi:10.1007/s00209-006-0072-5};
\href{https://arxiv.org/pdf/math/0505327v3}{arXiv:math/0505327v3}.

\bibitem{Horing26}
Andreas H\"oring,
\emph{Fano varieties with split tangent sheaf},
preprint, 2026.
\href{https://arxiv.org/abs/2602.15427v1}{arXiv:2602.15427v1}.

\bibitem{Ivashkovich1992}
S.~M.~Ivashkovich,
\emph{The Hartogs-type extension theorem for meromorphic maps into
compact K\"ahler manifolds},
Invent. Math. \textbf{109} (1992), 47--54.
\href{https://doi.org/10.1007/BF01232018}{doi:10.1007/BF01232018}.

\bibitem{Ivashkovich2008}
S.~Ivashkovich,
\emph{Extra extension properties of equidimensional holomorphic mappings:
results and open questions},
preprint, 2008.
\href{https://arxiv.org/abs/0810.4588v2}{arXiv:0810.4588v2}.

\bibitem{KM98}
J\'anos Koll\'ar and Shigefumi Mori,
\emph{Birational Geometry of Algebraic Varieties},
Cambridge Tracts in Mathematics, vol.~134,
Cambridge University Press, Cambridge, 1998.
\href{https://doi.org/10.1017/CBO9780511662560}{doi:10.1017/CBO9780511662560}.

\bibitem{MP09}
Jo\"el Merker and Egmont Porten,
\emph{The Hartogs extension theorem on $(n-1)$-complete complex spaces},
Journal f\"ur die reine und angewandte Mathematik
\textbf{637} (2009), 23--39.
\href{https://doi.org/10.1515/CRELLE.2009.088}{doi:10.1515/CRELLE.2009.088}.

\bibitem{NS24}
Yusuke Nakamura and Kohsuke Shibata,
\emph{A counterexample to the PIA conjecture for minimal log discrepancies},
Duke Mathematical Journal \textbf{175} (2026), no.~10, 1715--1736.
\href{https://arxiv.org/abs/2404.06164v4}{arXiv:2404.06164v4}.

\bibitem{OpenAIRCSplitting2026}
OpenAI,
\emph{Integrability of split tangent bundles on rationally connected manifolds},
OpenAI Math Release preprint
\href{https://github.com/openai/math/blob/main/preprints/Integrability-of-split-tangent-bundles-on-rationally-connected-manifolds-September-23-2026/main.pdf}{OAI:Integrability-of-split-tangent-bundles-on-rationally-connected-manifolds-September-23-2026},
2026.

\bibitem{PT13}
J.~V.~Pereira and F.~Touzet,
\emph{Foliations with vanishing Chern classes},
Bull. Braz. Math. Soc. (N.S.) \textbf{44} (2013), no.~4, 731--754.
\href{https://arxiv.org/abs/1210.5916v1}{arXiv:1210.5916v1}.

\bibitem{Serre56}
Jean-Pierre Serre,
\emph{G\'eom\'etrie alg\'ebrique et g\'eom\'etrie analytique},
Annales de l'Institut Fourier \textbf{6} (1956), 1--42.
\href{https://doi.org/10.5802/aif.59}{doi:10.5802/aif.59}.

\bibitem{Song14}
Jian Song,
\emph{Riemannian geometry of K\"ahler--Einstein currents},
preprint, 2014.
\href{https://arxiv.org/abs/1404.0445v1}{arXiv:1404.0445v1}.

\bibitem{Takayama03}
Shigeharu Takayama,
\emph{Local simple connectedness of resolutions of log-terminal singularities},
International Journal of Mathematics \textbf{14} (2003), no.~8, 825--836.
\href{https://doi.org/10.1142/S0129167X0300196X}{doi:10.1142/S0129167X0300196X}.

\bibitem{Yau93}
Shing-Tung Yau,
\emph{A splitting theorem and an algebraic geometric characterization of
locally Hermitian symmetric spaces},
Communications in Analysis and Geometry \textbf{1} (1993), 473--486.

\end{thebibliography}

\end{document}
